%% =====================================================================
%% P5 manuscript v3.1 (v1 = 9923d50, v2 = 1814564, v3 = ca0c350). v3: statistical errata E1-E6, coherent-joint
%% controls (P5-SYN-LEARN-01), calibration history moved to the appendix, body edited toward 8 pages.
%% v3.1 (round 5): wording corrections only; v4 (neural flow arms) is blocked on dependencies.
%%   TARGET_YEAR      = 2027  (intended venue: ICML 2027)
%%   TEMPLATE_YEAR    = 2026  (official ICML 2026 style, anonymous review mode;
%%                             no official ICML 2027 style was found on 2026-09-26)
%%   SUBMISSION_READY = false
%%   COMPILE_STATUS   = COMPILE_NOT_RUN (no LaTeX toolchain in the build environment)
%% The style files are NOT vendored in this repository. To build, place the
%% official files from https://media.icml.cc/Conferences/ICML2026/Styles/icml2026.zip
%% (icml2026.sty, icml2026.bst, algorithm.sty, algorithmic.sty, fancyhdr.sty)
%% next to this file and run: pdflatex main; bibtex main; pdflatex main; pdflatex main.
%% Margins, fonts and spacing of the style are not modified.
%% All numbers in the text come from generated/numbers.tex, which
%% scripts/make_paper_assets.py writes from committed raw results.
%% Unrun evidence is marked \todo{EXPERIMENT_ID: required evidence}.
%% =====================================================================
\documentclass{article}

\usepackage{microtype}
\usepackage{graphicx}
\usepackage{subcaption}
\usepackage{booktabs}
\usepackage{hyperref}

\newcommand{\theHalgorithm}{\arabic{algorithm}}

% Initial blind version for review:
\usepackage{icml2026}

\usepackage{amsmath}
\usepackage{amssymb}
\usepackage{mathtools}
\usepackage{amsthm}
\usepackage[capitalize,noabbrev]{cleveref}
\usepackage{pgfplots}
\pgfplotsset{compat=1.17}

\theoremstyle{plain}
\newtheorem{theorem}{Theorem}[section]
\newtheorem{proposition}[theorem]{Proposition}
\newtheorem{lemma}[theorem]{Lemma}
\newtheorem{corollary}[theorem]{Corollary}
\theoremstyle{definition}
\newtheorem{definition}[theorem]{Definition}
\newtheorem{assumption}[theorem]{Assumption}
\theoremstyle{remark}
\newtheorem{remark}[theorem]{Remark}

\newcommand{\todo}[1]{\textcolor{red}{\textbf{[TODO: #1]}}}
\newcommand{\notrun}{\textsc{not run}}
\newcommand{\ED}{\mathrm{ED}}
\newcommand{\E}{\mathbb{E}}
\newcommand{\qd}{q_{\mathrm{d}}}
\newcommand{\qs}{q_{\mathrm{s}}}

% AUTO-GENERATED by scripts/make_paper_assets.py from committed raw results. Do not edit by hand.
\newcommand{\numNullP}{0.026}
\newcommand{\numNullPjoint}{0.057}
\newcommand{\numMarkovNullP}{0.015}
\newcommand{\numRepRejT}{2/20}
\newcommand{\numRepRejJ}{1/20}
\newcommand{\numRepCPlo}{0.012}
\newcommand{\numRepCPhi}{0.317}
\newcommand{\numRepZsd}{1.33}
\newcommand{\numRepZsdLo}{1.01}
\newcommand{\numRepZsdHi}{1.95}
\newcommand{\numZflipSdMin}{0.985}
\newcommand{\numZflipSdMax}{1.014}
\newcommand{\numZflipQmin}{1.61}
\newcommand{\numZflipQmax}{1.69}
\newcommand{\numKurtMin}{4.3}
\newcommand{\numKurtMax}{7.6}
\newcommand{\numFloorM}{0.0082}
\newcommand{\numCRPSexact}{0.575}
\newcommand{\numCRPScollapse}{0.815}
\newcommand{\numCRPScollapsePct}{41.6}
\newcommand{\numCRPScollapseTheoryPct}{41.4}
\newcommand{\numCollIntD}{0.0227}
\newcommand{\numCollIntDiv}{0.73}
\newcommand{\numCollIntMSE}{1.54}
\newcommand{\numExactMSE}{2.02}
\newcommand{\numObsIntCRPS}{-0.180 [-0.229, -0.131]}
\newcommand{\numCorrPmin}{0.278}
\newcommand{\numCorrPmax}{0.623}
\newcommand{\numPredCellsT}{43}
\newcommand{\numPredCellsJ}{25}
\newcommand{\numPredZmaxT}{2.16}
\newcommand{\numPredZmaxJ}{2.60}
\newcommand{\numGapPrior}{0.00017}
\newcommand{\numGapCFGHalf}{0.031}
\newcommand{\numGapCFGOne}{0.155}
\newcommand{\numGapCFGTwo}{0.677}
\newcommand{\numGapCFGFour}{1.909}
\newcommand{\numGapEulerSixtyFour}{0.00017}
\newcommand{\numGapEulerFour}{0.00247}
\newcommand{\numBiasEulerSixteen}{0.0095}
\newcommand{\numBiasCFGFour}{0.077}
\newcommand{\numBiasCFGOne}{0.0058}
\newcommand{\numSEsolverMin}{0.0006}
\newcommand{\numSEsolverMax}{0.0009}
\newcommand{\numRuntimeVone}{80.9}
\newcommand{\numRuntimeEeight}{1.1}
\newcommand{\numEeightNullK}{9/200}
\newcommand{\numEeightNullLo}{0.024}
\newcommand{\numEeightNullHi}{0.077}
\newcommand{\numEeightShiftK}{42/50}
\newcommand{\numEeightShiftLo}{0.730}
\newcommand{\numEeightAdatasets}{55}
\newcommand{\numEeightBdatasets}{10}
\newcommand{\numEeightBminP}{0.0025}
\newcommand{\numCalibCPUmin}{24}
\newcommand{\numCollapseVsExactCRPS}{0.239 [0.205, 0.273]}
\newcommand{\numResolution}{0.00143}
\newcommand{\numOCvTwo}{0.121}
\newcommand{\numOCvTwoFail}{0.88}
\newcommand{\numOCvThree}{0.948}
\newcommand{\numOCvThreeGate}{0.945}
\newcommand{\numOCtolFailHundred}{0.56}
\newcommand{\numOCtolFailTwoHundred}{0.20}
\newcommand{\numOCprimFail}{0.019}
\newcommand{\numOCprimDetectTen}{0.964}
\newcommand{\numOCsecDetectTen}{0.42}
\newcommand{\numPres}{0.005}
\newcommand{\numPfamMax}{10}
\newcommand{\numCalibGate}{PASS}
\newcommand{\numCalibGOne}{pass}
\newcommand{\numCalibGTwo}{pass}
\newcommand{\numCalibGThree}{pass}
\newcommand{\numCalibPrimK}{16/400}
\newcommand{\numCalibPrimLo}{0.025}
\newcommand{\numCalibPrimHi}{0.060}
\newcommand{\numCalibPrimTwoLo}{0.023}
\newcommand{\numCalibPrimTwoHi}{0.064}
\newcommand{\numCalibWall}{19.9}
\newcommand{\numCalibCPU}{39.7}
\newcommand{\numCalibPowShift}{45/50}
\newcommand{\numCalibPowVar}{26/50}
\newcommand{\numCalibPowCollInt}{50/50}
\newcommand{\numCalibPowCorrJ}{50/50}
\newcommand{\numCalibNullJ}{5/100}
\newcommand{\numCalibNullCorrT}{5/100}
\newcommand{\numCalibNullChain}{4/100}
\newcommand{\numSynStatus}{RUN}
\newcommand{\numSynTrueP}{0.075}
\newcommand{\numSynTruePJ}{0.035}
\newcommand{\numSynFitP}{0.460}
\newcommand{\numSynFitPJ}{0.250}
\newcommand{\numSynHolmTrue}{0.15}
\newcommand{\numSynHolmFit}{0.46}
\newcommand{\numSynFitQ}{0.030 [0.025, 0.035]}
\newcommand{\numSynFitQJ}{0.021 [0.018, 0.024]}
\newcommand{\numSynTrueEJ}{2.08 [-1.67, 5.82]}
\newcommand{\numSynFitSdr}{1.056}
\newcommand{\numSynLLTrue}{-3.390}
\newcommand{\numSynLLFit}{-3.623}
\newcommand{\numSynWall}{22.5}
\newcommand{\numSynEMIter}{37}
\newcommand{\numFmStatus}{COMPLETE}
\newcommand{\numFmIndP}{0.005}
\newcommand{\numFmIndPJ}{0.005}
\newcommand{\numFmIndPThirtyTwo}{0.005}
\newcommand{\numFmIndHolm}{0.010}
\newcommand{\numFmIndD}{11.90 [7.19, 16.61]}
\newcommand{\numFmIndDJ}{11.83 [7.01, 16.65]}
\newcommand{\numFmIndQ}{1.91 [0.02, 3.81]}
\newcommand{\numFmIndQSeq}{9.84 [6.95, 12.74]}
\newcommand{\numFmIndQJ}{5.22 [2.95, 7.49]}
\newcommand{\numFmIndSdr}{0.980 [0.965, 0.995]}
\newcommand{\numFmIndShift}{-0.000 [-0.018, 0.017]}
\newcommand{\numFmIndMinDistinct}{64}
\newcommand{\numFmIndSolver}{1.10 [0.79, 1.41]}
\newcommand{\numFmIndSolverFlag}{yes}
\newcommand{\numFmIndSolverJ}{0.55 [0.31, 0.80]}
\newcommand{\numFmIndSolverQ}{-0.28 [-0.53, -0.03]}
\newcommand{\numFmIndParams}{34757}
\newcommand{\numFmIndUpdates}{20000}
\newcommand{\numFmIndTrainCPU}{32}
\newcommand{\numFmIndTrainWall}{33}
\newcommand{\numFmIndUnder}{-0.16}
\newcommand{\numFmIndUnderFlag}{no}
\newcommand{\numFmIndLabel}{EVALUABLE}
\newcommand{\numFmIndDevP}{0.005}
\newcommand{\numFmIndSampCPU}{14}
\newcommand{\numFmIndExposure}{z$\mid$x: 5000, y$\mid$x: 5000, z$\mid$x,y: 5000, y,z$\mid$x: 5000}
\newcommand{\numFmShP}{0.425}
\newcommand{\numFmShPJ}{0.035}
\newcommand{\numFmShPThirtyTwo}{0.295}
\newcommand{\numFmShHolm}{0.425}
\newcommand{\numFmShD}{0.24 [-2.75, 3.24]}
\newcommand{\numFmShDJ}{2.52 [0.28, 4.76]}
\newcommand{\numFmShQ}{5.51 [3.52, 7.51]}
\newcommand{\numFmShQSeq}{2.39 [0.40, 4.38]}
\newcommand{\numFmShQJ}{3.91 [2.59, 5.23]}
\newcommand{\numFmShSdr}{0.932 [0.919, 0.946]}
\newcommand{\numFmShShift}{-0.045 [-0.061, -0.028]}
\newcommand{\numFmShMinDistinct}{64}
\newcommand{\numFmShSolver}{-0.35 [-0.64, -0.06]}
\newcommand{\numFmShSolverFlag}{yes}
\newcommand{\numFmShSolverJ}{-0.25 [-0.42, -0.08]}
\newcommand{\numFmShSolverQ}{-1.48 [-1.76, -1.21]}
\newcommand{\numFmShParams}{34819}
\newcommand{\numFmShUpdates}{20000}
\newcommand{\numFmShTrainCPU}{47}
\newcommand{\numFmShTrainWall}{47}
\newcommand{\numFmShUnder}{1.04}
\newcommand{\numFmShUnderFlag}{no}
\newcommand{\numFmShLabel}{EVALUABLE}
\newcommand{\numFmShDevP}{0.585}
\newcommand{\numFmShSampCPU}{35}
\newcommand{\numFmShExposure}{z$\mid$x: 5052, y$\mid$x: 4985, z$\mid$x,y: 4953, y,z$\mid$x: 5010}
\newcommand{\numFmBetween}{-11.65 [-16.93, -6.38]}
\newcommand{\numFmBetweenBoot}{[-16.86, -6.30]}
\newcommand{\numFmBetweenN}{200}
\newcommand{\numFmBetweenJ}{-9.30 [-14.43, -4.18]}
\newcommand{\numFmBetweenThirtyTwo}{-10.20 [-15.32, -5.09]}
\newcommand{\numFmCPU}{220}
\newcommand{\numFmBudgetCPU}{238}
\newcommand{\numFmWall}{227}
\newcommand{\numFmTorch}{2.14.1{+}cpu}
\newcommand{\numRsevenStatus}{COMPLETE}
\newcommand{\numRsevenNreps}{4}
\newcommand{\numRsevenTheta}{-9.70 [-12.57, -6.83]}
\newcommand{\numRsevenThetaBoot}{[-12.47, -6.89]}
\newcommand{\numRsevenRepMeans}{-22.81, -11.18, -8.15, 3.35}
\newcommand{\numRsevenRepSd}{10.75}
\newcommand{\numRsevenReversal}{yes}
\newcommand{\numRsevenThetaJ}{-5.51 [-7.64, -3.37]}
\newcommand{\numRsevenThetaFiveTwelve}{-9.65 [-12.53, -6.77]}
\newcommand{\numRsevenTargetRejected}{3}
\newcommand{\numRsevenTargetRejectedInd}{3}
\newcommand{\numRsevenTargetRejectedSh}{0}
\newcommand{\numRsevenJointRejectedInd}{4}
\newcommand{\numRsevenJointRejectedSh}{2}
\newcommand{\numRsevenJointRawSh}{0.050, 0.010, 0.110, 0.010}
\newcommand{\numRsevenTargetRawSh}{0.085, 0.950, 0.900, 0.035}
\newcommand{\numRsevenTargetRawInd}{0.005, 0.005, 0.005, 0.230}
\newcommand{\numRsevenIndQdRange}{2.4\text{ to }4.6}
\newcommand{\numRsevenIndQsRange}{2.2\text{ to }14.5}
\newcommand{\numRsevenIndSolverFlags}{2}
\newcommand{\numRsevenIndTrainCPU}{26\text{ to }29}
\newcommand{\numRsevenIndUnderFlags}{1}
\newcommand{\numRsevenShQdRange}{2.8\text{ to }9.7}
\newcommand{\numRsevenShQsRange}{4.5\text{ to }13.0}
\newcommand{\numRsevenShSolverFlags}{2}
\newcommand{\numRsevenShTrainCPU}{37\text{ to }40}
\newcommand{\numRsevenShUnderFlags}{1}
\newcommand{\numRsevenCPU}{1494}
\newcommand{\numRsevenBudgetCPU}{1500}
\newcommand{\numRsevenWall}{1517}


\icmltitlerunning{Separating Sampling Error from Conditional Incompatibility}

\begin{document}

\twocolumn[
  \icmltitle{Separating Sampling Error from Conditional Incompatibility\\
             in Generative Models}

  \icmlsetsymbol{equal}{*}

  \begin{icmlauthorlist}
    \icmlauthor{Firstname1 Lastname1}{yyy}
  \end{icmlauthorlist}

  \icmlaffiliation{yyy}{Department of XXX, University of YYY, Location, Country}

  \icmlcorrespondingauthor{Firstname1 Lastname1}{first1.last1@xxx.edu}

  \icmlkeywords{conditional compatibility, marginalization consistency, any-to-any generation, two-sample testing, diffusion models, classifier-free guidance}

  \vskip 0.3in
]

\printAffiliationsAndNotice{}

\begin{abstract}
Any-to-any generators expose many conditionals of an implied joint, and users compose them by generating an intermediate modality and then the target.
If the conditionals are coherent, direct generation $q(z\mid x)$ and intermediate-then-target generation $\int q(z\mid x,y)\,q(y\mid x)\,dy$ agree in distribution. This marginalization-consistency condition is known, and it has been measured without sampling for models with explicit densities.
For generators that expose only samples, a measured direct-versus-sequential gap also contains Monte Carlo error, solver discretization, initialization mismatch, guidance and evaluation-design errors.
We give a conditional sampling design with the conditioning example as the independent unit, together with a within-example permutation test. We then separate these components on probes whose laws are known in closed form.
On a Gaussian probe, guidance with exact scores alone creates population gaps up to \numGapCFGFour{} (standardized), against a Monte Carlo floor of \numFloorM. Moderate solver error, by contrast, cancels between the two routes although the sampler is biased.
Our first calibration gate failed, and the next one was found defective before use. The test then passed a pre-registered calibration once.
On a Gaussian-mixture probe, a fitted normalized density, coherent by construction, was not flagged ($p=\numSynFitP$), although its distance to the truth is clearly nonzero.
Flow-matching learners were not run, and no real multimodal data are used.
\end{abstract}

% ---------------------------------------------------------------------
\section{Introduction}
\label{sec:intro}

Any-to-any generators such as UniDiffuser~\citep{bao2023unidiffuser}, Multi-modal Latent Diffusion~\citep{bounoua2023mld}, OmniFlow~\citep{li2025omniflow}, 4M~\citep{mizrahi2023fourm}, CoDi~\citep{tang2023codi} and MODUS~\citep{ye2026modus} expose conditionals for many subsets of modalities.
These conditionals are routinely composed.
UniDiffuser produces image variations by generating a caption and then an image from it, and runs blocked Gibbs sampling between modalities~\citep[Sec.~6.3]{bao2023unidiffuser}; 4M and MODUS generate chains of intermediate modalities~\citep{mizrahi2023fourm,ye2026modus}.
If all conditionals were induced by one joint distribution, the route would not matter in distribution: for an observation $x$, an intermediate $y$ and a target $z$,
\begin{equation}
  \qd(z\mid x) \;=\; \int q(z\mid x,y)\,q(y\mid x)\,dy \;=:\; \qs(z\mid x).
  \label{eq:identity}
\end{equation}
Equation~\eqref{eq:identity} is the law of total probability.
Whether a family of conditionals is compatible with some joint is a classical question~\citep{arnold1989compatible}.
For learned predictors, \citet{klotergens2026tabular} call \eqref{eq:identity} \emph{marginalization consistency} and show that tabular foundation models violate it.
None of this is our contribution.

What changes for diffusion and flow generators is \emph{how} \eqref{eq:identity} can be checked.
\citet{klotergens2026tabular} evaluate total variation from the models' predictive heads by exact sums or deterministic quadrature, with ``no sampling anywhere in the pipeline''.
A diffusion model exposes neither density, only a sampler.
A sampled discrepancy between the two routes then contains, besides any incompatibility of the learned conditionals, five other things:
(i)~Monte Carlo error;
(ii)~discretization error of the ODE or SDE solver;
(iii)~mismatch between the initial distribution and the noised data law;
(iv)~classifier-free guidance, which does not sample the intended conditional~\citep{chidambaram2024guidance,bradley2024cfg};
(v)~errors of evaluation design, such as feeding observed rather than generated intermediates, pairing samples by index, or counting generated samples as independent units.
Unless these are separated, a measured gap says little about the model.

This paper is about that separation, on synthetic probes only. We make three contributions.
\begin{itemize}
  \item \textbf{A conditional sampling design and its estimands} (\cref{sec:setup,sec:method}).
    \begin{itemize}
    \item The conditioning example is the independent unit. Intermediates are fresh for every sample, branches use independent randomness, and a permutation test relabels samples only within an example.
    \item We state what each estimand identifies: the target law, or eight projections of the joint.
    \item We keep apart three claims: non-detection, distributional correctness and the existence of a coherent joint.
    \end{itemize}
  \item \textbf{An error decomposition on closed-form probes} (\cref{sec:exp-numerics,sec:exp-calib}).
    \begin{itemize}
    \item With exact scores, the output law of a guided probability-flow ODE is available in closed form, so Monte Carlo error, discretization, prior mismatch and guidance can each be isolated.
    \item On this probe the test passed a pre-registered calibration once, after failures that we keep on record.
    \end{itemize}
  \item \textbf{Coherent-joint controls} (\cref{sec:exp-coherent}). Exact and fitted Gaussian mixtures are normalized densities, so their conditionals are coherent by construction while their fit quality varies. This separates coherence from correctness in a learned model.
\end{itemize}
The experiment that bears on the central question is not run: comparing separately learned conditionals with a shared conditional network (\cref{sec:exp-real}).
We do not claim as new:
\begin{itemize}
\item the identity~\eqref{eq:identity};
\item the consistency and compatibility definitions;
\item the fact that marginal-level consistency does not imply joint-level consistency, which is Proposition~2 of \citet{klotergens2026tabular};
\item guidance distortion;
\item permutation-test validity~\citep{ernst2004permutation,lehmann2005testing,phipson2010permutation};
\item energy statistics~\citep{szekely2013energy,baringhaus2004twosample}.
\end{itemize}

% ---------------------------------------------------------------------
\section{Related Work}
\label{sec:related}

\paragraph{Compatibility of conditionals.}
\citet{arnold1989compatible} characterize when conditionals admit a joint distribution. Compatible conditionals may still imply an improper joint~\citep{hobert1998functional}.
Pseudo-Gibbs sampling with incompatible conditionals has a scan-order-dependent limit~\citep{heckerman2000dependency,chen2015gibbs,liu2014imputation}.
We use only the definition, and ask how to measure one compatibility relation from samples.

\paragraph{Self-consistency of learned models.}
Masked language models produce conditionals that no single joint can explain~\citep{young2023mlm}. Any-order autoregressive models define the same distribution in redundant ways~\citep{shih2022aoarm}.
Generative marginalization models add a squared log-space penalty on single-step marginalization constraints for discrete data with explicit likelihoods~\citep{liu2024gmm}. This is the closest prior training-time consistency loss.
For diffusion language models, \citet{kim2026pathdependent} relate order dependence to incompatible denoising conditionals.
Closest to our estimand, \citet{klotergens2026tabular} define marginalization consistency (C1) and factorization consistency (C2) and prove that C2 implies C1 but not conversely. They measure both by total variation from the explicit predictive heads of tabular foundation models and find violations for every model and dataset.
We adopt their C1 as our target-level estimand. We add nothing to the definitions; our setting differs only in that the generator exposes samples.

\paragraph{Any-to-any generation and coherence.}
UniDiffuser trains one noise-prediction network with independent per-modality timesteps; a maximal timestep yields a marginal, which also serves as the unconditional model for guidance~\citep{bao2023unidiffuser}.
Multi-modal Latent Diffusion trains a masked multi-time score network on the latents of frozen unimodal autoencoders~\citep{bounoua2023mld}, and measures coherence with pre-trained classifiers, as in the multimodal VAE literature~\citep{wu2018mvae,shi2019mmvae,sutter2021mopoe,daunhawer2022limitations}.
Both are \emph{shared conditional networks}: one network serves many conditionals. That by itself does not make them the conditionals of one normalized joint.
\citet{chung2025anytoany} evaluate cyclic consistency and equivariance of any-to-any models with per-sample criteria.
None of these works tests \eqref{eq:identity} as an equality of distributions.

\paragraph{Guidance, solvers and testing.}
Classifier-free guidance~\citep{ho2022cfg} does not sample the guided target it is often described by~\citep{chidambaram2024guidance,bradley2024cfg}. We therefore treat it as a confound to be measured, not as model error.
We use the probability-flow ODE~\citep{song2021sde} with the step schedule of \citet{karras2022edm}.
For testing we use energy distances~\citep{szekely2013energy,baringhaus2004twosample}; kernel tests~\citep{gretton2012kernel} would serve equally well.
We use permutation p-values with the $+1$ correction~\citep{phipson2010permutation,ernst2004permutation} and the continuous ranked probability score (CRPS) as a strictly proper score~\citep{gneiting2007proper}.

% ---------------------------------------------------------------------
\section{Problem Setup}
\label{sec:setup}

A generator provides samplers for $\qd(z\mid x)$, $q(y\mid x)$ and $q(z\mid x,y)$, and possibly a direct joint sampler $\qd(y,z\mid x)$.
Here $x$ is observed, $y$ is an intermediate modality and $z$ is the target. Conditioning examples $x_1,\dots,x_N$ are drawn i.i.d.\ from a test split.

\begin{definition}[Target- and joint-level gaps]
For a distance $D$ between distributions, the \emph{target-level gap} is $\Delta_T(x) = D\big(\qd(\cdot\mid x),\,\qs(\cdot\mid x)\big)$ with $\qs$ as in \eqref{eq:identity}.
When a direct joint sampler exists, the \emph{joint-level gap} is $\Delta_J(x) = D\big(\qd(\cdot,\cdot\mid x),\,q(y\mid x)\,q(z\mid x,y)\big)$.
\end{definition}

$\Delta_T\equiv 0$ is C1 of \citet{klotergens2026tabular}.
$\Delta_J$ compares the sequential pair with a direct joint sampler. It is not their factorization consistency (C2), which compares the two chain-rule orders; we do not measure C2.

In practice we estimate a \emph{projected} joint gap. For a fixed set $U$ of eight unit directions,
\begin{equation}
  \Delta_J^{U}(x)=\frac{1}{|U|}\sum_{u\in U}D\big(u^{\top}_{\#}\qd(\cdot,\cdot\mid x),\,u^{\top}_{\#}\,q(y\mid x)q(z\mid x,y)\big),
  \label{eq:projected}
\end{equation}
where $u^\top_\#$ denotes the law of the projection $u^\top(y,z)$.
$\Delta_J^U=0$ is implied by joint equality but does not imply it: two joints can agree on finitely many projections and differ elsewhere.
Within the Gaussian family used below, eight directions identify the law; in general they do not.

\paragraph{Three different claims.}
We keep three statements apart; \cref{tab:claims} lists what would support each.
\emph{Non-detection}: a test at a given design ($N$, $M$, $B$) did not reject.
\emph{Distributional correctness}: a branch samples the law it is meant to sample, e.g.\ the true conditional.
\emph{Coherent joint existence}: a single joint distribution induces every conditional the model exposes.
Non-detection is bounded by power and resolution and does not establish $\Delta_T=0$. Even $\Delta_T=0$ does not establish correctness, as \cref{prop:collapse} and the common-mode solver bias of \cref{sec:exp-numerics} show. Equality of one relation is necessary for a coherent joint, but it is not sufficient.
This paper measures the first kind of statement for two relations, $\Delta_T$ and $\Delta_J^U$, together with quality metrics that bear on the second. The third holds by construction only for the coherent controls of \cref{rem:coherent}; it is measured for no model.

\begin{table}[t]
\centering
\caption{Three claims that a consistency check is often taken to support, and what supports each here.}
\label{tab:claims}
\scriptsize
\setlength{\tabcolsep}{3pt}
\begin{tabular}{p{0.24\columnwidth}p{0.36\columnwidth}p{0.3\columnwidth}}
\toprule
Claim & Evidence that would support it & Status in this paper \\
\midrule
Non-detection of $\Delta_T$, $\Delta_J^U$ & a calibrated, powered test that does not reject; reported with its design & synthetic probes and coherent controls; learned flows \notrun{} \\
Distributional correctness & proper scores or distances against a known truth & closed form on the probes \\
Coherent joint existence & all relations of the model, or a normalized density & by construction for controls only \\
\bottomrule
\end{tabular}
\end{table}

Neither check establishes a global joint for all conditionals of an any-to-any model. Each tests one relation.

\begin{proposition}[Target-level equality is compatibility of the triple]
\label{prop:triple}
There is a joint $r(y,z\mid x)$ whose $y\mid x$ marginal is $q(y\mid x)$, whose $z\mid x,y$ conditional is $q(z\mid x,y)$, and whose $z\mid x$ marginal is $\qd(z\mid x)$, if and only if $\qs(\cdot\mid x)=\qd(\cdot\mid x)$.
\end{proposition}

The proof is two lines (\cref{app:proofs}).
Because the first two properties force $r=q(y\mid x)\,q(z\mid x,y)$, target-level equality does not constrain the model's \emph{direct joint} sampler.

\begin{remark}[Target level does not determine the joint]
\label{rem:joint}
In a Gaussian model, replace the coefficient $b$ of $y-m_{y\mid x}(x)$ in $\E[z\mid x,y]$ by $cb$ and the residual variance by $\tau^2-c^2b^2v_y$, where $\tau^2=\mathrm{Var}(z\mid x)$ and $v_y=\mathrm{Var}(y\mid x)$.
Then $\qs(z\mid x)=p(z\mid x)$ for every $c$, while $\mathrm{Cov}(y,z\mid x)$ changes from $bv_y$ to $cbv_y$.
This is a continuous instance of the separation in Proposition~2 of \citet{klotergens2026tabular}, not a new result. We use it as a control that a target-level check must \emph{not} detect.
\end{remark}

\begin{proposition}[Collapse passes consistency]
\label{prop:collapse}
If every branch returns the conditional mean deterministically, then $\qd$ and $\qs$ are the same point mass and $\Delta_T=\Delta_J=0$.
For a Gaussian truth $\mathcal N(m,\sigma^2)$, the expected CRPS of that point mass is $\sigma\sqrt{2/\pi}$, against $\sigma/\sqrt{\pi}$ for the true conditional: $\sqrt2-1\approx\numCRPScollapseTheoryPct\%$ worse.
\end{proposition}

\begin{proposition}[Observed intermediates leak]
\label{prop:leak}
Let $(y^{\mathrm{obs}},z^{\mathrm{obs}})$ be drawn jointly given $x$, and let $S$ be a strictly proper, negatively oriented score.
Then $\E\,S\big(p(\cdot\mid x,y^{\mathrm{obs}}),z^{\mathrm{obs}}\big)\le\E\,S\big(p(\cdot\mid x),z^{\mathrm{obs}}\big)$.
\end{proposition}

This follows from propriety applied conditionally on $(x,y^{\mathrm{obs}})$ (\cref{app:proofs}).
So a ``sequential'' branch fed with the observed intermediate scores better than the direct branch while answering a different question. It must not be used to test \eqref{eq:identity}.
Chain generation $z\sim q(z\mid y)$, which drops $x$ at the second stage, agrees with $p(z\mid x)$ only if $z\perp x\mid y$. A mismatch there is expected and is not evidence of incompatibility.

\begin{remark}[Normalized joints are coherent controls]
\label{rem:coherent}
If all conditionals of a model are computed from one normalized density $r(x,y,z)$, then $\Delta_T=\Delta_J=0$ for every $x$ with $r(x)>0$ (\cref{prop:triple} with that $r$).
For a Gaussian mixture, the weights of $r(z\mid x,y)$ times $r(y\mid x)$ reduce to $\pi_k(x)\,\mathcal N_k(y\mid x)$, so the integral in \eqref{eq:identity} is available component-wise (\cref{app:proofs}).
Such a model can fit the data badly and still be exactly coherent. We call it a \textsc{joint\_coherent\_control}.
A network shared across conditionals, a \textsc{shared\_conditional\_net}, carries no such guarantee.
\end{remark}

% ---------------------------------------------------------------------
\section{Sampling Design, Tests and Error Decomposition}
\label{sec:method}

\subsection{Conditional sampling design}
\label{sec:design}

For each conditioning example $i$ we draw $A_i$, $M$ i.i.d.\ samples from the direct branch, and $C_i$, $M$ i.i.d.\ samples from the sequential branch.
Every sequential sample uses a fresh intermediate $\hat y_{ij}\sim q(y\mid x_i)$ followed by $z_{ij}\sim q(z\mid x_i,\hat y_{ij})$.
The first design (v1) also draws a second direct set $B_i$ to measure the Monte Carlo floor. For joint-level tests, a sample is the vector $(\hat y_{ij},z_{ij})$.
The assumptions behind every test below are:
\begin{assumption}
\label{ass:design}
(A1)~Examples are independent. (A2)~Within each example and branch, samples are i.i.d.\ given $x_i$. (A3)~The branches of an example are generated with independent randomness.
\end{assumption}
Three common shortcuts break these assumptions.
Generating several targets from one intermediate draw violates (A2): target samples then share a latent and are not exchangeable with direct samples.
Reusing initial noise across branches, i.e.\ common random numbers, violates (A3).
Pooling or relabelling samples across examples mixes different null distributions.
Generated samples, sign flips and permutations are never counted as independent units; only the $N$ examples are.

\subsection{Energy-distance estimators}
\label{sec:ed}

For samples $a_{1:n}$, $c_{1:m}$, the V-statistic energy distance is $\ED_V=\frac{2}{nm}\sum_{j,k}|a_j-c_k|-\frac{1}{n^2}\sum_{j,k}|a_j-a_k|-\frac{1}{m^2}\sum_{j,k}|c_j-c_k|$. The U-statistic $\ED_U$ divides the within-sample sums by $n(n-1)$ and $m(m-1)$ instead.
In one dimension both are computed exactly in $O((n+m)\log(n+m))$ time from sorted samples.
\begin{lemma}[standard]
\label{lem:vbias}
If $a$ is drawn i.i.d.\ from $P$ and $c$ i.i.d.\ from $Q$, then $\E\,\ED_U=D(P,Q)$ and $\E\,\ED_V=D(P,Q)+\E_P|X-X'|/n+\E_Q|Y-Y'|/m$, where $D$ is the population energy distance.
\end{lemma}
The v1 checker used the noise-floor-subtracted contrast $d_i=\ED_V(A_i,C_i)-\ED_V(A_i,B_i)$.
By \cref{lem:vbias}, $\E d_i=D+(\E|C-C'|-\E|B-B'|)/M$. This is exactly $0$ under $H_0$ and biased under alternatives whose spread differs from the reference; the bias is negative for under-dispersed $C$.
$\ED_U(A_i,C_i)$ is unbiased for $D$ and needs no third set.
We report effect sizes on a scale standardized by the training-split standard deviation of each coordinate.

\subsection{Two randomization tests}
\label{sec:tests}

\paragraph{Test S (v1): sign flips of $d_i$.}
Under $H_0$ and \cref{ass:design}, $B_i$ and $C_i$ are exchangeable given $A_i$, so $d_i$ is symmetric about $0$ and independent across $i$.
The one-sided sign-flip test of $\sum_i d_i$ therefore has level at most $\alpha$.
This is standard randomization logic. Each flip swaps the whole sets $B_i$ and $C_i$, so there are only two allocations per example and the noise of $B_i$ enters $d_i$.

\paragraph{Test P (revised): relabelling within examples.}
Under $H_0$ and \cref{ass:design}, the $2M$ pooled samples of example $i$ are exchangeable.
We reassign branch labels uniformly over all $\binom{2M}{M}$ allocations independently in every example, never across examples, and relabel vectors as units.
The pre-specified statistic is $T=\frac1N\sum_i\ED_V(a_i,c_i)$ in the upper tail.
The Monte Carlo p-value is $(1+\#\{b:T_b\ge T_{\mathrm{obs}}\})/(B+1)$, with ties counted as extreme, so it is never zero~\citep{phipson2010permutation}. On small designs the product group is enumerated exactly.
\begin{proposition}[V and U give the same permutation test]
\label{prop:vu}
With equal group sizes, the permutation p-values of $T$ built from $\ED_V$ and from $\ED_U$ coincide.
\end{proposition}
Given the pooled multiset of an example, the within- and cross-sums add to a constant, so both statistics are increasing affine functions of the cross-sum with slopes depending only on $M$ (\cref{app:proofs}).
The choice between V and U thus matters for effect sizes, not for the test.

\paragraph{What Test P can and cannot establish.}
For the target level, Test~P has level at most $\alpha$ under $H_0^T$: $\qd(z\mid x_i)=\qs(z\mid x_i)$ for all $i$.
For the joint level, the pooled vectors are exchangeable only under equality of the \emph{full} joint laws, so the level guarantee holds under that null. Its statistic depends only on the projections in $U$, so it has power only against differences visible in them (\cref{eq:projected}).
A rejection therefore shows that the joint laws differ. A non-rejection says nothing about directions outside $U$.

\paragraph{Resolution and multiplicity.}
With $B$ permutations the p-value lies on a grid of step $1/(B+1)$; for $B=199$ the smallest value is $\numPres$.
Several conditions together make $P(p\le\alpha)\le\alpha$ under $H_0$~\citep{phipson2010permutation}:
\begin{itemize}
\item the observed allocation is counted ($+1$) and ties are counted as extreme;
\item the $B$ allocations are drawn i.i.d.\ and uniformly from the within-example relabelling group;
\item the statistic is fixed in advance.
\end{itemize}
Our implementation meets all three.
Equality $P(p\le\alpha)=\lfloor\alpha(B+1)\rfloor/(B+1)$ needs in addition that the observed and permuted statistics have no ties almost surely. That fails in general: the group is finite, $\ED_V$ is invariant under complementing an allocation, and collapsed samples tie. We therefore claim level $\le\alpha$, not exact size.
A Bonferroni family of $m$ tests can reject at all only if $m\le\alpha(B+1)=\numPfamMax$.

\subsection{Error decomposition on a Gaussian probe}
\label{sec:decomp}

Energy distances do not add, so we define components operationally: each is the gap that remains when all other sources are switched off in a probe where every law is computable.
\begin{itemize}
\item Monte Carlo error vanishes as $M\to\infty$ (\cref{lem:vbias}).
\item Discretization vanishes as the number of solver steps $N_s\to\infty$.
\item Prior mismatch is the residual gap of the exact flow started from $\mathcal N(0,\sigma_{\max}^2)$ rather than from the noised data law.
\item Guidance vanishes at $w=0$.
\item Model incompatibility is whatever remains in a learned model.
\end{itemize}

\begin{proposition}[Guided PF-ODE on Gaussian conditionals is affine]
\label{prop:affine}
Let the target be $\mathcal N(m_c,v_c)$ and the unconditional law $\mathcal N(m_u,v_u)$. Use exact noised scores in the variance-exploding parameterization, guided score $(1+w)s_c-ws_u$, initialization $x_0\sim\mathcal N(0,\sigma_{\max}^2)$, and Euler or Heun steps on any schedule.
Then the output is $x_N=g\,x_0+h_c m_c+h_u m_u$, with $(g,h_c,h_u)$ depending only on $(v_c,v_u,w)$ and the schedule.
\end{proposition}

The velocity $\sigma\big[(1+w)\frac{x-m_c}{v_c+\sigma^2}-w\frac{x-m_u}{v_u+\sigma^2}\big]$ is affine in $(x,m_c,m_u)$, and so is each Euler or Heun step.
Composing the two stages of the sequential branch keeps the output Gaussian, so both branch laws are available per example in closed form. We verify this against the sampler to $10^{-8}$ in unit tests; that is a numerical check, and the proof is the affinity argument itself.
The continuous-time limit is also closed-form: the exact flow keeps $(x-m)/\sqrt{v+\sigma^2}$ invariant.

\subsection{Quality metrics}
Consistency is always reported next to quality, because \cref{prop:collapse} shows that consistency can be satisfied by a useless generator.
\begin{itemize}
\item \textbf{Quality to truth.} Where the truth is known, we report the population energy distance between the model's conditional and the true one. For a one-dimensional law, half of it is the CRPS divergence~\citep{gneiting2007proper}.
\item \textbf{Proper score.} Otherwise we use the fair ensemble CRPS against held-out data.
\item \textbf{Diversity.} The sample standard deviation over the true conditional one.
\end{itemize}
These answer a different question from the compatibility test. Their detection rates are not comparable to its power.
Same-index paired MSE between branches is $2\,\mathrm{Var}$ for independent exact samplers and $0$ for collapsed ones, so it is not a compatibility metric.

% ---------------------------------------------------------------------
\section{Experiments}
\label{sec:exp}

\subsection{Probes and protocol}
\label{sec:exp-protocol}
Two synthetic probes are used; neither is real multimodal data.
\begin{itemize}
\item \textbf{Gaussian probe.} Three scalars $(x,y,z)$ with a non-Markov and a Markov correlation structure (\cref{app:probe}). Its laws are closed-form for exact samplers and for the PF-ODE.
\item \textbf{Mixture probe.} A fixed three-component Gaussian mixture over $(x,y,z)$ with multimodal, non-Gaussian conditionals (\cref{app:mixture}).
\end{itemize}
In both, splits are independent draws. The train split is used for fitting and standardization, the dev split only for pipeline smoke runs, and the test split ($N=200$ conditioning examples) for reported results.
Model, data and sampling seeds are separate. Every configuration was committed before its test run (\cref{app:prereg}).
Closed-form curves were computed after the corresponding runs and are exploratory.
The v1 Gaussian-probe runs used Test~S with $M=256$ and 999 flips; the later runs use Test~P with $N=200$, $M=64$ and $B=199$.
The v1 controls, including an exact null that the single test rejected, are reported in \cref{app:controls}.

% Figure layout (data come from paper/generated/*.dat, written by scripts/make_paper_assets.py).
% Palette: two categorical slots validated with the dataviz validator (light mode, all checks pass).
\definecolor{seriesA}{HTML}{2A78D6}
\definecolor{seriesB}{HTML}{EB6834}
\definecolor{refgray}{HTML}{52514E}
\pgfplotsset{p5axis/.style={
  width=\linewidth, height=0.82\linewidth,
  tick label style={font=\scriptsize}, label style={font=\scriptsize},
  title style={font=\scriptsize}, legend style={font=\tiny, draw=none, fill=none},
  grid=major, major grid style={draw=black!10}, axis line style={draw=black!45}}}
\begin{figure*}[t]
\centering
\begin{minipage}[t]{0.32\textwidth}
\begin{tikzpicture}
\begin{loglogaxis}[p5axis, title={(a) Monte Carlo noise floor},
  xlabel={samples per branch $M$}, ylabel={$\overline{\mathrm{ED}_V(A,B)}$ (standardised)},
  log basis x=2, xtick={32,64,128,256,512}, xticklabels={32,64,128,256,512},
  legend pos=north east]
\addplot[seriesA, thick, dashed] table[x=M, y=theory] {generated/fig_mc.dat};
\addlegendentry{$2\,\mathbb{E}|X-X'|/M$}
\addplot[seriesA, only marks, mark=*, mark size=1.6pt,
  error bars/.cd, y dir=both, y explicit]
  table[x=M, y=floor, y error plus expr=\thisrow{hi}-\thisrow{floor},
        y error minus expr=\thisrow{floor}-\thisrow{lo}] {generated/fig_mc.dat};
\addlegendentry{observed, 95\% CI}
\end{loglogaxis}
\end{tikzpicture}
\end{minipage}\hfill
\begin{minipage}[t]{0.32\textwidth}
\begin{tikzpicture}
\begin{loglogaxis}[p5axis, title={(b) Euler PF-ODE, exact scores, $w=0$},
  xlabel={solver steps $N$}, ylabel={population $D$ (standardised)},
  log basis x=2, xtick={2,4,8,16,32,64,128,256}, xticklabels={2,4,8,16,32,64,128,256},
  legend pos=north east]
\addplot[seriesB, thick] table[x=steps, y=euler_bias] {generated/fig_solver_closed.dat};
\addlegendentry{direct vs.\ exact law}
\addplot[seriesA, thick] table[x=steps, y=euler_seq] {generated/fig_solver_closed.dat};
\addlegendentry{direct vs.\ sequential}
\addplot[refgray, dashed] table[x=x, y=y] {generated/fig_res_steps.dat};
\addlegendentry{detection resolution}
\end{loglogaxis}
\end{tikzpicture}
\end{minipage}\hfill
\begin{minipage}[t]{0.32\textwidth}
\begin{tikzpicture}
\begin{semilogyaxis}[p5axis, title={(c) CFG, Euler 64 steps, exact scores},
  xlabel={guidance weight $w$}, ylabel={population $D$ / observed $\bar d$},
  xmin=0, xmax=4, legend pos=south east]
\addplot[seriesA, thick] table[x=w, y=seq] {generated/fig_cfg_closed.dat};
\addlegendentry{direct vs.\ sequential}
\addplot[seriesB, thick] table[x=w, y=bias] {generated/fig_cfg_closed.dat};
\addlegendentry{direct vs.\ exact law}
\addplot[seriesA, only marks, mark=*, mark size=1.6pt,
  error bars/.cd, y dir=both, y explicit]
  table[x=x, y=d, y error plus expr=\thisrow{hi}-\thisrow{d},
        y error minus expr=\thisrow{d}-\thisrow{lo}] {generated/fig_cfg_obs.dat};
\addlegendentry{observed $\bar d$, 95\% CI}
\addplot[refgray, dashed] table[x=x, y=y] {generated/fig_res_w.dat};
\addlegendentry{detection resolution}
\end{semilogyaxis}
\end{tikzpicture}
\end{minipage}
\caption{Separating numerical components on the Gaussian probe. (a) The Monte Carlo floor of the V-statistic energy distance between two independent sets follows $2\,\mathbb{E}|X-X'|/M$ (P5-E5, $N=200$). (b) With exact scores and no guidance, the closed-form gap between the direct and sequential PF-ODE laws falls below the one-sided 5\% detection resolution of the v1 design ($N=100$, $M=256$) from about 6 steps on, and plateaus at the prior-mismatch value, although the sampler remains biased relative to the exact law. (c) Guidance alone creates direct-vs-sequential gaps far above the resolution; the observed means (P5-E7) match the closed form. Closed-form curves are computed from the affine output law (\cref{prop:affine}); no Monte Carlo is involved in them.}
\label{fig:numerics}
\end{figure*}


\subsection{Where sampled gaps come from}
\label{sec:exp-numerics}


\paragraph{Monte Carlo.}
The floor $\overline{\ED_V(A,B)}$ follows $2\,\E|X-X'|/M$ from $M=32$ to $512$ (\cref{fig:numerics}a).
A $0.1\tau$ mean shift is detected at every $M$ in that range with $N=200$. At $M=512$, $\bar d$ is $5.1\times10^{-3}$ against a population value of $5.6\times10^{-3}$.

\paragraph{Discretization and prior mismatch.}
With exact scores and $w=0$, the Euler direct-vs-sequential gap is $\numGapEulerFour$ at $4$ steps and is detected ($p=0.014$, v1 Test~S).
From about $6$ steps on it is below the one-sided detection resolution of this design, $\numResolution$ (\cref{fig:numerics}b; \cref{tab:numerics} in \cref{app:numerics}). It plateaus at $\numGapEulerSixtyFour$ at $64$ steps, which equals the closed-form continuous-time limit, $\numGapPrior$. The plateau is therefore prior mismatch at $\sigma_{\max}=80$, not discretization.
Over the same range the direct sampler itself stays biased relative to the exact law. At $16$ steps its gap to the exact law is $\numBiasEulerSixteen$, which is detected, and its sd ratio is $0.83$.
Solver error is thus largely common to the two branches and cancels in the comparison: \emph{passing the consistency check does not certify a correct sampler}, and ``not detected'' is not ``zero''.
Heun with $4$ steps is unstable on this schedule (sd ratio $4.35$).

\paragraph{Guidance.}
Guidance alone produces population gaps of $\numGapCFGHalf$, $\numGapCFGOne$, $\numGapCFGTwo$ and $\numGapCFGFour$ at $w=0.5,1,2,4$, against a floor of about $0.008$ (\cref{fig:numerics}c). The observed means agree with the closed form.
The guided direct sampler stays close to the exact law: $D=\numBiasCFGOne$ at $w=1$ and $\numBiasCFGFour$ at $w=4$. Because $\sqrt{D}$ is a metric~\citep{szekely2013energy}, most of the direct-vs-sequential gap must come from the sequential branch, which applies guidance at both stages.
Under our pre-registered decision rule, a direct-vs-sequential gap measured with guidance on is not attributed to the model.
This agrees with prior analyses of guidance~\citep{chidambaram2024guidance,bradley2024cfg}. The magnitudes are specific to this probe.

\subsection{A calibrated permutation test on the Gaussian probe}
\label{sec:exp-calib}
Test~P passed one pre-registered calibration (\cref{tab:calib}; $N=200$, $M=64$, $B=199$; \numCalibWall{} minutes on two processes).
\begin{itemize}
\item \textbf{Primary exact null.} \numCalibPrimK{} rejections, two-sided 95\% Clopper--Pearson interval [\numCalibPrimTwoLo, \numCalibPrimTwoHi]. The gate's one-sided 95\% upper bound was \numCalibPrimHi{} $\le 0.10$.
\item \textbf{Other null cells.} Three cells showed no excess at the Bonferroni level: the exact null at the joint level, a correlation flip that preserves the target law, and a Markov chain.
\item \textbf{Alternatives.} Four fixed alternatives were all powered: \numCalibPowShift, \numCalibPowVar, \numCalibPowCollInt{} and \numCalibPowCorrJ.
\item \textbf{Quality control.} A both-branch collapse is a compatibility null and a quality failure. It was never rejected, and the quality metrics flagged it in every replicate.
\end{itemize}
This is engineering evidence for this implementation and design.
The level-$\alpha$ guarantee comes from exchangeability and a fixed statistic, not from this run. It transfers to a new model only if \cref{ass:design} holds for that model's sampling pipeline.
The secondary null cells detect a true size of $0.10$ with probability only $\numOCsecDetectTen$.
The route to this result included a failed gate and a gate withdrawn before use, both kept on record (\cref{app:calibhistory}).
\begin{enumerate}
\item The first checker's pre-registered gate failed on its single exact-null test ($p=\numNullP$). An audit traced this to the gate design rather than the estimator.
\item The next gate, audited before use, would have failed a test with null rejection rate $0.05$ with probability $\numOCvTwoFail$.
\end{enumerate}

\begin{table}[t]
\centering
\caption{Calibration of Test~P (single run). Rejection counts over independent replicates, with the one-sided 95\% bounds used by the gate (together an equal-tailed two-sided 90\% interval) and the two-sided 95\% Clopper--Pearson interval.}
\label{tab:calib}
\scriptsize
\setlength{\tabcolsep}{2pt}
% AUTO-GENERATED by scripts/make_paper_assets.py from committed raw results. Do not edit by hand.
\begin{tabular}{llrlll}
\toprule
Cell & family & rejections & 1-s.\ 95\% bounds & 2-s.\ 95\% CP & outcome \\
\midrule
Exact sequential & null validity & 16/400 & 0.025, 0.060 & [0.023, 0.064] & within tolerance, no excess \\
Exact sequential (joint) & null validity & 5/100 & 0.020, 0.102 & [0.016, 0.113] & no excess \\
Corr.\ scale $-1$ & null validity & 5/100 & 0.020, 0.102 & [0.016, 0.113] & no excess \\
Chain $z\mid y$ (Markov) & null validity & 4/100 & 0.014, 0.089 & [0.011, 0.099] & no excess \\
Mean shift $0.1\tau$ & power & 45/50 & 0.801, 0.960 & [0.782, 0.967] & POWERED \\
Variance ratio 1.25 & power & 26/50 & 0.395, 0.643 & [0.374, 0.663] & POWERED \\
Collapse, intermediate only & power & 50/50 & 0.942, 1.000 & [0.929, 1.000] & POWERED \\
Corr.\ scale $-1$ (joint) & power & 50/50 & 0.942, 1.000 & [0.929, 1.000] & POWERED \\
Collapse, both branches & quality control & 0/10 & 0.000, 0.259 & [0.000, 0.308] & not rejected; quality flagged \\
Collapse, both branches (joint) & quality control & 0/10 & 0.000, 0.259 & [0.000, 0.308] & not rejected; quality flagged \\
\bottomrule
\end{tabular}

\end{table}

\subsection{Coherent-joint controls on the mixture probe}
\label{sec:exp-coherent}
Two models are compared, each a normalized density whose conditionals are coherent by \cref{rem:coherent}.
\begin{itemize}
\item \textsc{true\_joint\_control}: the exact source mixture.
\item \textsc{fitted\_joint\_control}: a two-component mixture fitted by EM on the training split (\numSynEMIter{} iterations). It is misspecified by design, since the source has three components.
\end{itemize}
All conditionals of each model come from its own density and are sampled exactly, so no solver is involved.
The primary endpoint is the target test. The projected-joint test is secondary.
\Cref{tab:synlearn} reports the single pre-registered run (\numSynWall\,s).

\begin{table}[t]
\centering
\caption{Coherent-joint controls on the mixture probe (Test~P, $N=200$, $M=64$, $B=199$). $\widehat\Delta_T$ is the mean unbiased energy distance between the two branches. $D(\text{model},\text{truth})$ is the closed-form energy distance between the model's and the true conditional of $z$ given $x$, averaged over test examples. Both are standardized.}
\label{tab:synlearn}
\scriptsize
\setlength{\tabcolsep}{2pt}
% AUTO-GENERATED by scripts/make_paper_assets.py from committed raw results. Do not edit by hand.
\begin{tabular}{lrrrrr}
\toprule
Model & $p_T$ & $p_J$ & $\widehat{\Delta}_T\times10^3$ [95\% CI] & $D(\text{model},\text{truth})\times10^3$ [95\% CI] & test log-lik. \\
\midrule
True joint (exact mixture) & 0.075 & 0.035 & 2.71 [-1.38, 6.80] & 0.00 [0.00, 0.00] & -3.390 \\
Fitted joint (EM, $K=2$) & 0.460 & 0.250 & 0.21 [-3.03, 3.46] & 30.10 [25.26, 34.95] & -3.623 \\
Independent conditional FM & \multicolumn{5}{l}{NOT RUN (no tensor library; installation not approved)} \\
Shared conditional FM & \multicolumn{5}{l}{NOT RUN (no tensor library; installation not approved)} \\
\bottomrule
\end{tabular}

\end{table}

\paragraph{Coherence.}
Neither control is flagged by the primary target test: $p=\numSynTrueP$ and $\numSynFitP$, Holm-adjusted over the primary family $\numSynHolmTrue$ and $\numSynHolmFit$. The gap estimates contain $0$.
The exact mixture's secondary projected-joint test gave $p=\numSynTruePJ$, with gap estimate $\numSynTrueEJ\times10^{-3}$.
Its population gap is exactly zero, so this is a false positive of a level-$0.05$ test, one of four tests run. We report it as observed; it is not evidence of incompatibility.
The coherence of both controls follows from their construction (\cref{rem:coherent}), not from these $p$-values; non-rejection is only consistent with it.

\paragraph{Correctness.}
The fitted control is coherent yet wrong.
\begin{itemize}
\item Its conditional of $z$ given $x$ lies at energy distance $\numSynFitQ$ from the truth, and its projected joint at $\numSynFitQJ$.
\item Its test log-likelihood is $\numSynLLFit$, against $\numSynLLTrue$ for the source.
\item Its conditional spread is too large by a factor of $\numSynFitSdr$.
\end{itemize}
A compatibility test cannot see this by construction; only the quality metrics can.
\Cref{tab:decision} summarizes what the run does and does not decide.

\begin{table}[t]
\centering
\caption{Decision table for the mixture-probe pilot.}
\label{tab:decision}
\scriptsize
\setlength{\tabcolsep}{3pt}
\begin{tabular}{p{0.30\columnwidth}p{0.62\columnwidth}}
\toprule
Question & Status \\
\midrule
Coherence vs.\ quality & controls: coherent by construction, not flagged at the primary endpoint; the fitted control's quality is clearly worse \\
Separate vs.\ shared learners & \notrun{} (no tensor library; installation not approved) \\
Gap from solver refinement & not applicable to the controls (exact sampling); \notrun{} for learners \\
Non-detection vs.\ equality, correctness, global joint & for the controls equality and a global joint hold by construction, correctness fails for the fitted one; for the neural flow learners (\notrun) none is established \\
\bottomrule
\end{tabular}
\end{table}

\subsection{Neural flow learners (not run)}
\label{sec:exp-real}
The experiment that bears on the central question is designed and coded but not run. It compares independently trained conditional flow-matching models with one shared conditional network, both trained on the mixture probe.
The environment has no tensor library, and installing one was not approved.
\Cref{app:mixture} fixes the architecture, update budget, conditioning-pattern exposure and the two solver levels, so that data cannot influence them.
A static review before any execution found that the shared network could not tell an absent coordinate from a target one; its input now carries a target mask.
If the two solver levels disagree or training is insufficient, the outcome is labeled inconclusive, not incompatible.
\todo{P5-REAL-01: $\widehat\Delta_T$, $\widehat\Delta_J^U$, quality and cost for independent vs.\ shared conditional flow matching on the mixture probe}.
A consistency regularizer, and any real multimodal data, are outside this paper (\todo{P5-REAL-02: real paired multimodal data}).

% ---------------------------------------------------------------------
\section{Limitations}
\label{sec:limits}
\begin{itemize}
\item \textbf{Data and models.} All evidence is synthetic. The only learned model is an EM-fitted mixture; no neural flow or diffusion generator has been evaluated. Nothing here shows that such models are, or are not, incompatible.
\item \textbf{Closed-form decomposition.} It relies on Gaussian conditionals and exact scores. With learned scores, discretization, prior mismatch and model error interact and must be separated empirically.
\item \textbf{Joint test.} It compares eight projections and cannot certify joint equality.
\item \textbf{Calibration.} It ran once, at one design. Its secondary null cells screen only for gross anti-conservativeness.
\item \textbf{Validity scope.} The level guarantee requires fresh intermediates and independent branch randomness. Pipelines that reuse one intermediate for several targets need another design.
\item \textbf{Mixture pilot.} It used one data seed, one model seed and one fitted control.
\item \textbf{Guidance convention.} The unconditional model is the target marginal, following \citet{bao2023unidiffuser}; other conventions give other gaps.
\item \textbf{Scope of the discrepancy.} Energy distance is one of several possible discrepancies. Factorization consistency (C2) is not measured.
\item \textbf{Build.} The manuscript has not been compiled in our environment.
\end{itemize}

% ---------------------------------------------------------------------
\section{Conclusion}
\label{sec:conclusion}
For samplers without densities, checking whether direct and intermediate-then-target generation agree is mostly a matter of accounting for other sources of error.
On probes whose laws are known, guidance can dominate the gap, and moderate solver error cancels between routes even when the sampler is wrong.
A permutation test with the conditioning example as the unit passed a pre-registered calibration once, after one failed gate and one gate design withdrawn before use.
Coherent-by-construction controls are not flagged even when their fit is poor. Consistency is therefore neither correctness nor evidence of a global joint.
Whether learned any-to-any generators are incompatible beyond these effects remains open.

% ---------------------------------------------------------------------
\section*{Impact Statement}
This paper presents work whose goal is to advance the field of Machine Learning. There are many potential societal consequences of our work, none of which we feel must be specifically highlighted here.
The work concerns evaluation methodology. A test that confuses sampler or guidance artifacts with model incoherence could lead practitioners to wrong conclusions about generated multimodal content. We therefore report our calibration failure and all unrun evidence explicitly.
No human or personal data were used.

\bibliography{references}
\bibliographystyle{icml2026}

% =====================================================================
\newpage
\appendix
\onecolumn

\section{Proofs}
\label{app:proofs}

\paragraph{\Cref{prop:triple}.}
($\Leftarrow$) Take $r(y,z\mid x)=q(y\mid x)\,q(z\mid x,y)$. Its $y\mid x$ marginal and $z\mid x,y$ conditional are the given ones, and its $z\mid x$ marginal is $\qs=\qd$.
($\Rightarrow$) Any $r$ with the stated $y\mid x$ marginal and $z\mid x,y$ conditional equals $q(y\mid x)\,q(z\mid x,y)$ wherever $q(y\mid x)>0$. Its $z\mid x$ marginal is therefore $\qs$, which must equal $\qd$. \qed

\paragraph{\Cref{prop:collapse}.}
Deterministic branches that return $\E[z\mid x]$, respectively $\E[z\mid x,\E[y\mid x]]$, coincide for Gaussian conditionals, so both laws are $\delta_{\E[z\mid x]}$.
The CRPS of a point mass at $m$ against $Z\sim\mathcal N(m,\sigma^2)$ is $\E|Z-m|=\sigma\sqrt{2/\pi}$.
The expected CRPS of $\mathcal N(m,\sigma^2)$ against its own draw is $\E|Z-m|-\tfrac12\E|Z-Z'|=\sigma\sqrt{2/\pi}-\sigma/\sqrt{\pi}=\sigma/\sqrt\pi$. The ratio of the two is $\sqrt2$. \qed

\paragraph{\Cref{prop:leak}.}
Conditionally on $(x,y^{\mathrm{obs}})$ we have $z^{\mathrm{obs}}\sim p(\cdot\mid x,y^{\mathrm{obs}})$. By propriety, $\E[S(p(\cdot\mid x,y^{\mathrm{obs}}),z^{\mathrm{obs}})\mid x,y^{\mathrm{obs}}]\le\E[S(Q,z^{\mathrm{obs}})\mid x,y^{\mathrm{obs}}]$ for every $Q$, in particular for $Q=p(\cdot\mid x)$. Averaging over $(x,y^{\mathrm{obs}})$ gives the claim. \qed

\paragraph{\Cref{lem:vbias}.}
The $n^2$ within-sample terms include $n$ zero diagonal terms, so $\E\frac{1}{n^2}\sum_{j,k}|a_j-a_k|=(1-\frac1n)\E|X-X'|$, and likewise for $c$. Hence $\E\,\ED_V=2\E|X-Y|-(1-\frac1n)\E|X-X'|-(1-\frac1m)\E|Y-Y'|=D+\frac1n\E|X-X'|+\frac1m\E|Y-Y'|$. The U-statistic removes the diagonal terms. This result is standard~\citep{szekely2013energy}; we include the derivation because the v1 contrast depends on it. \qed

\paragraph{Validity of Test~S.}
Under $H_0$ and \cref{ass:design}, $(A_i,B_i,C_i)$ and $(A_i,C_i,B_i)$ have the same distribution, so $d_i\overset{d}{=}-d_i$. The $d_i$ are independent across $i$, so conditional on $(|d_i|)_i$ the signs are i.i.d.\ uniform. The Monte Carlo sign-flip p-value with the $+1$ correction is then valid~\citep{phipson2010permutation}. \qed

\paragraph{\Cref{rem:coherent} for Gaussian mixtures.}
Let $r=\sum_k\pi_k\mathcal N_k(x,y,z)$, and write $\pi_k(x)=\pi_k\mathcal N_k(x)/r(x)$ and $w_k(x,y)=\pi_k\mathcal N_k(x,y)/r(x,y)$.
Then $r(y\mid x)\,w_k(x,y)=\frac{r(x,y)}{r(x)}\cdot\frac{\pi_k\mathcal N_k(x,y)}{r(x,y)}=\pi_k(x)\,\mathcal N_k(y\mid x)$.
Hence $\int r(z\mid x,y)\,r(y\mid x)\,dy=\sum_k\pi_k(x)\int\mathcal N_k(z\mid x,y)\,\mathcal N_k(y\mid x)\,dy=\sum_k\pi_k(x)\,\mathcal N_k(z\mid x)=r(z\mid x)$. \qed
A quadrature check of this identity is in the unit tests; it is a numerical check, not part of the proof.

\paragraph{\Cref{prop:vu}.}
Fix example $i$ and its pooled multiset. Let $S_{aa}$, $S_{cc}$ and $S_{ac}$ be the within-$a$, within-$c$ and cross sums of absolute differences over unordered pairs. Then $S_{aa}+S_{cc}+S_{ac}=S_{\mathrm{pool}}$ is invariant under relabelling.
With $n=m=M$, $\ED_V=\frac{2S_{ac}}{M^2}-\frac{2(S_{\mathrm{pool}}-S_{ac})}{M^2}$ and $\ED_U=\frac{2S_{ac}}{M^2}-\frac{2(S_{\mathrm{pool}}-S_{ac})}{M(M-1)}$.
Both are increasing affine functions of $S_{ac}$, with slopes that depend only on $M$ and intercepts that are invariant under relabelling.
Summing over examples, both versions of $T$ are the same increasing affine function of $\sum_i S_{ac,i}$ up to invariant constants, so they order all allocations identically. \qed

\paragraph{\Cref{prop:affine}.}
With $s_c=-(x-m_c)/(v_c+\sigma^2)$ and $s_u=-(x-m_u)/(v_u+\sigma^2)$, the velocity $dx/d\sigma=-\sigma[(1+w)s_c-ws_u]$ equals $\sigma a(\sigma)x-\sigma[b_c(\sigma)m_c+b_u(\sigma)m_u]$, with $a=\frac{1+w}{v_c+\sigma^2}-\frac{w}{v_u+\sigma^2}$, $b_c=\frac{1+w}{v_c+\sigma^2}$ and $b_u=-\frac{w}{v_u+\sigma^2}$.
An Euler step $x\mapsto x+h\,v(x,\sigma)$ and a Heun step are compositions of affine maps in $(x,m_c,m_u)$ with coefficients independent of $(x,m_c,m_u)$. Induction over the steps gives the claim. \qed

\section{Probe Details}
\label{app:probe}
\begin{itemize}
\item \textbf{Conditionals} come from the Schur complement $T\mid G=g\sim\mathcal N(\mu_T+\Sigma_{TG}\Sigma_{GG}^{-1}(g-\mu_G),\,\Sigma_{TT}-\Sigma_{TG}\Sigma_{GG}^{-1}\Sigma_{GT})$.
\item \textbf{Perturbations} act on the second stage $q(z\mid x,y)$:
  \begin{itemize}
  \item a mean shift adds $\delta\tau$ to the intercept;
  \item a variance ratio $r$ sets the residual variance to $r\tau^2-b^2v_y$;
  \item a correlation scale $c$ follows \cref{rem:joint};
  \item a stage-1 shift adds $\delta\sqrt{v_y}$ to $\E[y\mid x]$.
  \end{itemize}
\item \textbf{Collapse} sets the sampling temperature to $0$.
\item \textbf{Sampler.} The PF-ODE uses the schedule of \citet{karras2022edm} with $\sigma_{\min}=0.002$, $\sigma_{\max}=80$, $\rho=7$, the final step to $\sigma=0$, and initialization $\mathcal N(0,\sigma_{\max}^2)$.
\item \textbf{Guidance.} The unconditional model for both stages is the marginal of the stage target.
\item \textbf{Scale.} All discrepancies are computed after dividing $y$ and $z$ by their train-split standard deviations.
\item \textbf{Sample reuse.} Reference sets are shared across conditions within a run (common random numbers). The Monte Carlo experiment uses nested prefixes across $M$.
\end{itemize}

\section{Estimator Versus Closed Form for All v1 Cells}
\label{app:reanalysis}\label{sec:exp-estimator}

For every exact-Gaussian and PF-ODE condition, the closed form of \cref{lem:vbias,prop:affine} predicts $\E\bar d$ from the laws of the branches alone.
Across \numPredCellsT{} non-degenerate target-level cells, the observed means lie within $|z|\le\numPredZmaxT$ of the prediction, and across \numPredCellsJ{} joint-level cells within $|z|\le\numPredZmaxJ$ (\cref{app:reanalysis}).
The two target-level cells with $|z|>2$ are the exact nulls on the two joints, which share their random numbers; see \cref{sec:calib}.
The spread term of \cref{lem:vbias} is visible where it should be. For variance ratio $0.8$ the population $D$ is $3.3\times10^{-3}$, while $\E\bar d=2.9\times10^{-3}$.

\Cref{tab:reanalysis} lists, for every v1 condition, the population gap, the closed-form expectation of the v1 contrast and the observed mean with its 95\% CI. It was computed after the results were seen and is exploratory. $z$ is not defined for the collapsed condition, whose contrasts are identically $0$.
\begin{table}[h]
\centering
\caption{Closed-form expectation versus observation, all v1 cells (values $\times10^3$).}
\label{tab:reanalysis}
\scriptsize
% AUTO-GENERATED by scripts/make_paper_assets.py from committed raw results. Do not edit by hand.
\begin{tabular}{llrrrrrrr}
\toprule
Exp. & Condition & $M$ & $D_{\mathrm{pop}}$ & $\mathbb{E}[d]$ & $\bar d$ [95\% CI] & $z$ & $\mathbb{E}[d_J]$ & $\bar d_J$ \\
 & & & \multicolumn{6}{c}{(all $\times10^3$)} \\
\midrule
E1-NULL & \texttt{exact\_sequential} & 256 & 0.000 & 0.000 & 1.23 [0.03, 2.43] & 2.00 & -0.000 & 0.64 \\
E2-CHECKER & \texttt{mean\_shift\_0.05} & 256 & 1.412 & 1.412 & 2.03 [0.54, 3.53] & 0.81 & 0.821 & 1.12 \\
E2-CHECKER & \texttt{mean\_shift\_0.1} & 256 & 5.646 & 5.646 & 4.39 [2.80, 5.98] & -1.55 & 3.284 & 3.46 \\
E2-CHECKER & \texttt{mean\_shift\_0.2} & 256 & 22.558 & 22.558 & 22.22 [19.15, 25.29] & -0.22 & 13.118 & 13.31 \\
E2-CHECKER & \texttt{mean\_shift\_0.5} & 256 & 139.768 & 139.768 & 136.98 [130.19, 143.78] & -0.80 & 81.201 & 81.13 \\
E2-CHECKER & \texttt{var\_ratio\_0.8} & 256 & 3.321 & 2.855 & 2.13 [1.00, 3.25] & -1.27 & 2.049 & 2.74 \\
E2-CHECKER & \texttt{var\_ratio\_1.1} & 256 & 0.657 & 0.872 & 0.54 [-0.59, 1.66] & -0.59 & 0.542 & 0.97 \\
E2-CHECKER & \texttt{var\_ratio\_1.25} & 256 & 3.713 & 4.234 & 4.81 [3.37, 6.26] & 0.79 & 2.592 & 2.91 \\
E2-CHECKER & \texttt{var\_ratio\_1.5} & 256 & 12.793 & 13.784 & 13.31 [11.41, 15.21] & -0.49 & 8.216 & 8.46 \\
E2-CHECKER & \texttt{corr\_scale\_-1} & 256 & 0.000 & 0.000 & 0.07 [-0.98, 1.12] & 0.14 & 65.473 & 66.25 \\
E2-CHECKER & \texttt{corr\_scale\_0} & 256 & 0.000 & 0.000 & -0.16 [-1.37, 1.06] & -0.25 & 18.050 & 17.71 \\
E2-CHECKER & \texttt{corr\_scale\_0.5} & 256 & 0.000 & 0.000 & 0.36 [-0.76, 1.49] & 0.63 & 5.615 & 6.62 \\
E2-CHECKER & \texttt{stage1\_mean\_shift\_0.5} & 256 & 64.684 & 64.684 & 64.42 [59.73, 69.11] & -0.11 & 78.882 & 80.35 \\
E3-COLLAPSE & \texttt{collapse\_both} & 256 & 0.000 & 0.000 & 0.00 [0.00, 0.00] & NA & 0.000 & 0.00 \\
E3-COLLAPSE & \texttt{collapse\_intermediate} & 256 & 22.789 & 21.618 & 22.75 [20.97, 24.52] & 1.25 & 109.412 & 111.60 \\
E4-INTERMEDIATE & \texttt{observed\_intermediate} & 256 & 323.306 & 322.135 & 320.78 [269.99, 371.57] & -0.05 & -- & -- \\
E4-INTERMEDIATE & \texttt{chain\_drop\_x} & 256 & 40.017 & 39.742 & 39.28 [31.13, 47.44] & -0.11 & 23.491 & 23.41 \\
E4-MARKOV & \texttt{exact\_sequential} & 256 & -0.000 & -0.000 & 1.20 [0.11, 2.29] & 2.16 & -0.000 & 0.58 \\
E4-MARKOV & \texttt{chain\_drop\_x} & 256 & -0.000 & -0.000 & -0.18 [-1.08, 0.71] & -0.40 & -0.000 & 0.19 \\
E6-SOLVER & \texttt{euler\_2\_seq} & 256 & 0.751 & 0.778 & 0.78 [0.74, 0.83] & 0.23 & -- & -- \\
E6-SOLVER & \texttt{euler\_4\_seq} & 256 & 2.472 & 2.186 & 1.21 [0.17, 2.25] & -1.84 & -- & -- \\
E6-SOLVER & \texttt{euler\_8\_seq} & 256 & 0.263 & 0.343 & -0.47 [-1.59, 0.64] & -1.43 & -- & -- \\
E6-SOLVER & \texttt{euler\_16\_seq} & 256 & 0.162 & 0.189 & 0.20 [-1.17, 1.56] & 0.01 & -- & -- \\
E6-SOLVER & \texttt{euler\_32\_seq} & 256 & 0.162 & 0.171 & -0.39 [-1.95, 1.17] & -0.71 & -- & -- \\
E6-SOLVER & \texttt{euler\_64\_seq} & 256 & 0.166 & 0.166 & 1.03 [-0.68, 2.73] & 0.99 & -- & -- \\
E6-SOLVER & \texttt{heun\_4\_seq} & 256 & 41.779 & 38.414 & 40.88 [31.69, 50.08] & 0.53 & -- & -- \\
E6-SOLVER & \texttt{heun\_16\_seq} & 256 & 0.196 & 0.159 & 0.72 [-1.04, 2.49] & 0.63 & -- & -- \\
E6-SOLVER & \texttt{euler\_4\_direct\_vs\_exact} & 256 & 80.997 & 78.910 & 79.14 [75.66, 82.61] & 0.13 & -- & -- \\
E6-SOLVER & \texttt{euler\_16\_direct\_vs\_exact} & 256 & 9.461 & 8.699 & 9.56 [7.50, 11.61] & 0.82 & -- & -- \\
E6-SOLVER & \texttt{euler\_64\_direct\_vs\_exact} & 256 & 0.894 & 0.699 & 0.85 [-0.43, 2.13] & 0.23 & -- & -- \\
E7-CFG & \texttt{cfg\_0.5\_seq} & 256 & 30.620 & 30.792 & 31.16 [20.20, 42.12] & 0.07 & -- & -- \\
E7-CFG & \texttt{cfg\_1\_seq} & 256 & 155.058 & 155.270 & 160.18 [117.73, 202.64] & 0.23 & -- & -- \\
E7-CFG & \texttt{cfg\_2\_seq} & 256 & 676.600 & 676.538 & 674.62 [513.78, 835.45] & -0.02 & -- & -- \\
E7-CFG & \texttt{cfg\_4\_seq} & 256 & 1908.864 & 1907.677 & 1893.88 [1509.16, 2278.60] & -0.07 & -- & -- \\
E7-CFG & \texttt{cfg\_1\_direct\_vs\_exact} & 256 & 5.760 & 5.544 & 5.71 [2.80, 8.63] & 0.11 & -- & -- \\
E7-CFG & \texttt{cfg\_4\_direct\_vs\_exact} & 256 & 76.689 & 76.409 & 84.87 [60.90, 108.84] & 0.69 & -- & -- \\
E5-MC & \texttt{null\_m32} & 32 & 0.000 & 0.000 & 5.25 [-4.72, 15.22] & 1.03 & -0.000 & 0.80 \\
E5-MC & \texttt{null\_m64} & 64 & 0.000 & 0.000 & 0.24 [-4.49, 4.98] & 0.10 & -0.000 & -0.52 \\
E5-MC & \texttt{null\_m128} & 128 & 0.000 & 0.000 & 0.61 [-1.51, 2.74] & 0.56 & -0.000 & 1.03 \\
E5-MC & \texttt{null\_m512} & 512 & 0.000 & 0.000 & -0.11 [-0.65, 0.43] & -0.40 & -0.000 & -0.10 \\
E5-MC & \texttt{mean\_shift\_0.1\_m32} & 32 & 5.646 & 5.646 & 14.98 [3.86, 26.10] & 1.64 & 3.284 & 5.72 \\
E5-MC & \texttt{mean\_shift\_0.1\_m64} & 64 & 5.646 & 5.646 & 8.00 [2.71, 13.29] & 0.87 & 3.284 & 6.06 \\
E5-MC & \texttt{mean\_shift\_0.1\_m128} & 128 & 5.646 & 5.646 & 6.11 [3.08, 9.13] & 0.30 & 3.284 & 4.25 \\
E5-MC & \texttt{mean\_shift\_0.1\_m512} & 512 & 5.646 & 5.646 & 5.14 [4.12, 6.17] & -0.96 & 3.284 & 2.92 \\
\bottomrule
\end{tabular}

\end{table}

\section{Normal Reference Check for the Studentized Contrast}
\label{app:znull}
\begin{table}[h]
\centering
\caption{Sign-flip distribution of the studentized mean of $d_i$ given the observed $|d_i|$ (4999 flips) for the seven v1 null cells.}
\label{tab:znull}
\small
% AUTO-GENERATED by scripts/make_paper_assets.py from committed raw results. Do not edit by hand.
\begin{tabular}{llrrrrr}
\toprule
Exp. & Condition & skew$(d)$ & kurt$(d)$ & $t_{\mathrm{obs}}$ & sd of $t^\ast$ & $q_{0.95}(t^\ast)$ \\
\midrule
E1-NULL & \texttt{exact\_sequential} & 0.46 & 5.02 & 2.00 & 0.985 & 1.612 \\
E4-MARKOV & \texttt{exact\_sequential} & 0.30 & 4.28 & 2.16 & 1.014 & 1.644 \\
E4-MARKOV & \texttt{chain\_drop\_x} & -0.13 & 5.00 & -0.40 & 1.014 & 1.634 \\
E5-MC & \texttt{null\_m32} & -0.78 & 7.58 & 1.03 & 1.005 & 1.686 \\
E5-MC & \texttt{null\_m64} & -0.04 & 5.58 & 0.10 & 1.000 & 1.641 \\
E5-MC & \texttt{null\_m128} & 0.25 & 4.47 & 0.56 & 1.002 & 1.620 \\
E5-MC & \texttt{null\_m512} & 0.08 & 5.87 & -0.40 & 0.995 & 1.637 \\
\bottomrule
\end{tabular}

\end{table}

\section{Numerical Components: Closed Form Versus Observation}
\label{app:numerics}
\begin{table}[h]
\centering
\caption{Numerical components with exact scores ($N=100$, $M=256$, Test~S). $D_{\mathrm{pop}}$ is the closed-form population gap between the direct and sequential sampler laws, and $\E[d]$ is the closed-form expectation of the v1 contrast. All values $\times10^3$, standardized.}
\label{tab:numerics}
\scriptsize
\setlength{\tabcolsep}{3pt}
% AUTO-GENERATED by scripts/make_paper_assets.py from committed raw results. Do not edit by hand.
\begin{tabular}{lrrrr}
\toprule
Sampler (both branches) & $D_{\mathrm{pop}}\times10^3$ & $\mathbb{E}[d]\times10^3$ & $\bar d\times10^3$ [95\% CI] & $p_T$ \\
\midrule
Euler, 2 steps & 0.751 & 0.778 & 0.78 [0.74, 0.83] & 0.001 \\
Euler, 4 steps & 2.472 & 2.186 & 1.21 [0.17, 2.25] & 0.014 \\
Euler, 8 steps & 0.263 & 0.343 & -0.47 [-1.59, 0.64] & 0.804 \\
Euler, 64 steps & 0.166 & 0.166 & 1.03 [-0.68, 2.73] & 0.121 \\
Heun, 4 steps & 41.779 & 38.414 & 40.88 [31.69, 50.08] & 0.001 \\
Heun, 16 steps & 0.196 & 0.159 & 0.72 [-1.04, 2.49] & 0.206 \\
Euler 64, $w=0.5$ & 30.620 & 30.792 & 31.16 [20.20, 42.12] & 0.001 \\
Euler 64, $w=1$ & 155.058 & 155.270 & 160.18 [117.73, 202.64] & 0.001 \\
Euler 64, $w=4$ & 1908.864 & 1907.677 & 1893.88 [1509.16, 2278.60] & 0.001 \\
Exact flow, $N\to\infty$ & 0.173 & -- & NOT RUN & -- \\
\bottomrule
\end{tabular}

\end{table}

\section{Small Exact Verification (P5-E8-EXACT)}
\label{app:e8}
\begin{table}[h]
\centering
\caption{P5-E8-EXACT, Test~P with exact enumeration of the product permutation group ($N=3$ examples, $M=3$, $|G|=20^3$), nominal $\alpha=0.05$, one-sided 95\% Clopper--Pearson bounds on the rejection rate over independent datasets. ``Excess flagged'' means the lower bound exceeds $\alpha$; ``powered'' means the lower bound exceeds $\alpha$ under an alternative. Two negative-control rules are evaluated with the same criteria.}
\label{tab:e8}
\scriptsize
\setlength{\tabcolsep}{2.5pt}
% AUTO-GENERATED by scripts/make_paper_assets.py from committed raw results. Do not edit by hand.
\begin{tabular}{llrrrl}
\toprule
Cell & Role & rejections & 1-s.\ 95\% lower & 1-s.\ 95\% upper & Verdict \\
\midrule
Exact sequential, target & null & 9/200 & 0.024 & 0.077 & no excess flagged \\
Corr.\ scale $-1$, target & null & 4/100 & 0.014 & 0.089 & no excess flagged \\
Exact sequential, joint & null & 1/50 & 0.001 & 0.091 & no excess flagged \\
Mean shift $2\tau$, target & alternative & 42/50 & 0.730 & 0.918 & powered \\
Variance ratio 4, target & alternative & 5/50 & 0.040 & 0.199 & not powered \\
Corr.\ scale $-1$, joint & alternative & 4/50 & 0.028 & 0.174 & not powered \\
Collapse intermediate, target & alternative & 2/50 & 0.007 & 0.121 & not powered \\
\midrule
Always-accept rule & null / alt. & 0/200 & -- & -- & no excess flagged / not powered \\
Always-reject rule & null & 200/200 & -- & -- & excess flagged \\
\bottomrule
\end{tabular}

\end{table}

\begin{enumerate}
\item \textbf{Conditional size.} For all \numEeightAdatasets{} small datasets, including tied data from collapsed samplers, full enumeration confirms that the fraction of allocations with $p\le\alpha$ is at most $\alpha$ for $\alpha\in\{0.01,0.05,0.1,0.2,0.5\}$.
\item \textbf{Monte Carlo agreement.} For \numEeightBdatasets{} datasets, the Monte Carlo p-value with $B=1999$ agrees with the exact one within tolerance and is never below $1/(B+1)$; the smallest observed is $\numEeightBminP$.
\item \textbf{Size and power (\cref{tab:e8}).} The exact null rejects \numEeightNullK{} (one-sided bounds [\numEeightNullLo, \numEeightNullHi]), and the target-level correlation null and the joint-level null show no excess. A $2\tau$ mean shift is powered (\numEeightShiftK). Variance, correlation and collapse alternatives are \emph{not} powered at this tiny size, as recorded without adjustment.
\item \textbf{Negative controls.} The always-accept rule fails the power criterion, and the always-reject rule is flagged for excess.
\end{enumerate}

\section{Quality Metrics in the Calibration Run}
\label{app:calibq}
\begin{table}[h]
\centering
\caption{Quality metrics in the calibration run: mean CRPS excess over the true conditional ($\times10^3$) for each branch, oracle sd ratio of the sequential branch, and the share of replicates whose 95\% CI for the sequential CRPS excess lies above $0$.}
\label{tab:calibq}
\scriptsize
\setlength{\tabcolsep}{3pt}
% AUTO-GENERATED by scripts/make_paper_assets.py from committed raw results. Do not edit by hand.
\begin{tabular}{lrrrr}
\toprule
Cell & $\overline{\Delta\mathrm{CRPS}}$ seq.\ & $\overline{\Delta\mathrm{CRPS}}$ direct & sd ratio seq.\ & reps with CI $>0$ \\
\midrule
Exact sequential & 0.39 & 0.03 & 0.996 & 3\% \\
Exact sequential (joint) & -1.00 & 0.79 & 0.996 & 3\% \\
Corr.\ scale $-1$ & 0.24 & 0.57 & 0.996 & 3\% \\
Chain $z\mid y$ (Markov) & -0.47 & -0.24 & 0.997 & 3\% \\
Mean shift $0.1\tau$ & 3.90 & 0.19 & 0.995 & 12\% \\
Variance ratio 1.25 & 1.61 & -0.65 & 1.113 & 14\% \\
Collapse, intermediate only & 12.20 & 0.13 & 0.732 & 44\% \\
Corr.\ scale $-1$ (joint) & 1.52 & 1.79 & 0.995 & 6\% \\
Collapse, both branches & 235.32 & 235.32 & 0.000 & 100\% \\
Collapse, both branches (joint) & 235.11 & 235.11 & 0.000 & 100\% \\
\bottomrule
\end{tabular}

\end{table}

\section{v1 Controls with Test S}
\label{app:controls}\label{sec:exp-controls}
These are the first checker's pre-registered engineering controls on the Gaussian probe. They used Test~S with $N=200$ and $M=256$. The single exact-null test rejected. That failure is discussed in \cref{app:calibhistory}.
\begin{table}[h]
\centering
\caption{Controls on the non-Markov probe (last row: Markov), Test~S, $N=200$ examples, $M=256$. $\bar d$ and $\bar d_J$ are the mean target- and joint-level contrasts $\times10^3$. $\Delta$CRPS is the mean paired difference in CRPS between the compared and the reference branch $\times10^3$; $^\dagger$ marks the difference against the exact direct sampler, re-aggregated from raw. The sd ratio is the oracle sample sd over the true conditional sd. Every pre-registered detect / no-detect expectation in this table was met except for the exact null in the first row. The exact null on the Markov joint (not shown, $p=\numMarkovNullP$) also rejected; it reused the same noise (\cref{sec:calib}). Rows for $0.05\tau$ and variance ratio $1.1$ were pre-registered as descriptive.}
\label{tab:controls}
\scriptsize
% AUTO-GENERATED by scripts/make_paper_assets.py from committed raw results. Do not edit by hand.
\begin{tabular}{lrrrrrr}
\toprule
Condition & $\bar d\times10^3$ [95\% CI] & $p_T$ & $\bar d_J\times10^3$ & $p_J$ & $\Delta\mathrm{CRPS}\times10^3$ [95\% CI] & sd ratio \\
\midrule
Exact sequential (null) & 1.23 [0.03, 2.43] & 0.026 & 0.64 & 0.057 & 3.2 [-5.2, 11.6] & 1.00 \\
Mean shift $0.05\tau$ & 2.03 [0.54, 3.53] & 0.006 & 1.12 & 0.020 & 4.5 [-4.6, 13.7] & 1.00 \\
Mean shift $0.1\tau$ & 4.39 [2.80, 5.98] & 0.001 & 3.46 & 0.001 & 7.2 [-2.8, 17.1] & 1.00 \\
Mean shift $0.5\tau$ & 136.98 [130.19, 143.78] & 0.001 & 81.13 & 0.001 & 61.8 [21.2, 102.5] & 1.01 \\
Variance ratio 0.8 & 2.13 [1.00, 3.25] & 0.002 & 2.74 & 0.001 & 6.1 [-1.2, 13.4] & 0.89 \\
Variance ratio 1.1 & 0.54 [-0.59, 1.66] & 0.162 & 0.97 & 0.017 & 1.2 [-6.0, 8.3] & 1.05 \\
Variance ratio 1.5 & 13.31 [11.41, 15.21] & 0.001 & 8.46 & 0.001 & 12.3 [2.3, 22.4] & 1.22 \\
Corr.\ scale $-1$ (target exact) & 0.07 [-0.98, 1.12] & 0.463 & 66.25 & 0.001 & 0.7 [-7.4, 8.8] & 1.00 \\
Corr.\ scale $0.5$ (target exact) & 0.36 [-0.76, 1.49] & 0.278 & 6.62 & 0.001 & 1.4 [-6.6, 9.4] & 1.00 \\
Collapse, both branches & 0.00 [0.00, 0.00] & 1.000 & 0.00 & 1.000 & 239.2 [205.3, 273.0]$^\dagger$ & 0.00 \\
Collapse, intermediate only & 22.75 [20.97, 24.52] & 0.001 & 111.60 & 0.001 & 17.4 [5.5, 29.4] & 0.73 \\
Observed intermediate & 320.78 [269.99, 371.57] & 0.001 & -- & -- & -179.7 [-228.5, -130.8] & 0.73 \\
Chain $z\mid y$ (non-Markov) & 39.28 [31.13, 47.44] & 0.001 & 23.41 & 0.001 & 22.6 [0.4, 44.7] & 0.93 \\
Chain $z\mid y$ (Markov) & -0.18 [-1.08, 0.71] & 0.667 & 0.19 & 0.291 & -1.5 [-8.5, 5.4] & 1.00 \\
\bottomrule
\end{tabular}

\end{table}


\Cref{tab:controls} summarizes the pre-registered controls; these are engineering checks of the checker, not scientific findings.
Mean shifts of at least $0.1\tau$ and variance ratios of $0.8$, $1.25$ and $1.5$ are detected at the target level. A $0.05\tau$ shift is also detected ($p=0.006$; descriptive).
A variance ratio of $1.1$ is not detected at the target level ($p=0.162$) but is detected at the joint level.
The correlation breaks of \cref{rem:joint} leave the target marginal exact. They are not detected at the target level ($p=\numCorrPmin$ to $\numCorrPmax$) and are detected at the joint level ($p=0.001$).
Collapsing both branches gives $\bar d=0$ exactly and $p=1$, so the consistency check passes. Its CRPS is worse than the exact sampler's by \numCollapseVsExactCRPS{} ($\numCRPScollapse$ vs.\ $\numCRPSexact$, \numCRPScollapsePct\% against the theoretical \numCRPScollapseTheoryPct\%), and its paired-index MSE is $0$ against \numExactMSE{} for exact samplers.
Collapsing only the intermediate is detected ($\bar d=\numCollIntD$) and yields sd ratio \numCollIntDiv; the invalid paired-index MSE ``improves'' to \numCollIntMSE.
Feeding the observed intermediate is detected, and its CRPS is \emph{better} than the direct branch's by \numObsIntCRPS, the leakage predicted by \cref{prop:leak}.
Chain generation is detected under the non-Markov joint and not under the Markov joint.


\section{Calibration Development Record}
\label{app:calibhistory}\label{sec:calib}
This appendix records the route to the calibrated test of \cref{sec:exp-calib}, including both failures. Statistical wording follows the errata in \cref{app:errata}.


\paragraph{The failure.}
The v1 pre-registration required that the exact null not be detected and that a 20-replicate false-positive rate be compatible with $0.05$.
The single null test rejected at $p=\numNullP$ (joint level $p=\numNullPjoint$), so the gate \emph{failed}, and by our stop rule the v1 checker was not cleared for use on models.
We keep this verdict.
The replicates rejected \numRepRejT{} times at the target level (Clopper--Pearson 95\% [\numRepCPlo, \numRepCPhi]) and \numRepRejJ{} at the joint level. The replicate $z$ statistics had standard deviation \numRepZsd{} ($\chi^2$ 95\% CI [\numRepZsdLo, \numRepZsdHi]).

\paragraph{The audit.}
Three findings separate the gate from the estimator.
\begin{enumerate}
\item The gate was mis-specified. An exact $\alpha$-level test rejects a single null with probability $\alpha$. ``Compatible with $0.05$'' over 20 replicates carries little information. The gate had no power requirement, so an always-accept rule would have passed it.
\item The v1 seeds omitted the joint and experiment names. The exact-null rejection on the Markov joint ($p=\numMarkovNullP$) reused the same noise as the non-Markov one and is not independent evidence.
\item Test~S is exact under \cref{ass:design} (\cref{sec:tests}), and the normal reference used for the reported $z$ is adequate. Given the observed $|d_i|$ of each of the seven v1 null cells, the sign-flip distribution of the studentized mean has standard deviation \numZflipSdMin--\numZflipSdMax{} and 95\% quantile \numZflipQmin--\numZflipQmax{} (normal: $1.645$), although $d_i$ has kurtosis \numKurtMin--\numKurtMax{} (\cref{app:znull}).
\end{enumerate}
None of this shows that the checker is calibrated. It shows that the evidence does not demonstrate invalidity, and that the gate could not have established validity.

\paragraph{The revised test, checked exactly.}\label{sec:e8summary}
We replaced Test~S by Test~P, introduced v2 seeds that hash every design factor, and pre-registered a small exact verification (P5-E8-EXACT) and a separate realistic calibration (P5-E8-CALIB) before running either.

On the smallest designs the checks behave as they must (\cref{app:e8}).
Full enumeration confirms conditional size $\le\alpha$ for all \numEeightAdatasets{} datasets, including ties. Monte Carlo and exact p-values agree and are never $0$. The exact null rejects \numEeightNullK{} times.
A $2\tau$ shift is powered, but the variance, correlation and collapse alternatives are not powered at $N=M=3$. The always-accept and always-reject rules fail.
These are implementation checks of a textbook principle at $|G|\le 8000$, not a calibration of the checker at the size used in \cref{sec:exp-controls}.

\paragraph{Auditing the calibration contract before running it.}
The first realistic calibration design (P5-E8-CALIB v2) was never run.
A closed-form audit of its null and alternative cells, and of its gate, found three problems.
\begin{enumerate}
\item Its tolerance rule, a one-sided Clopper--Pearson upper bound $\le0.10$, would fail a test with null rejection rate $0.05$ (the least favourable rate for a level-$0.05$ test) with probability $\numOCtolFailHundred$ in a cell with $R=100$ and $\numOCtolFailTwoHundred$ with $R=200$. At that rate, a test would have passed all v2 null cells with probability only $\numOCvTwo$.
\item The seed key of the small exact verification omitted the test level, so its target- and joint-level cells shared data.
\item There was no quality-control family, so the check could not show that a compatibility null which is a quality failure is handled correctly.
\end{enumerate}

\Cref{tab:contract} states what each cell changes, in closed form.
Three alternatives change the target law, the projected joint law and the distance to the truth. The correlation flip changes only the joint, so it is a null for the target test and an alternative for the joint test. Both-branch collapse changes neither branch law relative to the other, so it is a compatibility null, yet it is far from the truth.

The corrected design (v3) was fixed before any calibration result. It uses $N=200$, $M=64$ and $B=199$, with seeds that include the level, and three separate families.
\begin{itemize}
\item \textbf{Null validity:} four cells.
  \begin{itemize}
  \item The primary cell, exact sequential at the target level, has $R=400$ and a tolerance rule. At a null rejection rate of $0.05$ it fails with probability $\numOCprimFail$; a test of size $0.10$ is caught with probability $\numOCprimDetectTen$.
  \item The other three cells have $R=100$ and are checked only by excess flags at the Bonferroni level over the four cells. They catch a size of $0.10$ with probability only $\numOCsecDetectTen$.
  \end{itemize}
\item \textbf{Power:} four fixed alternatives, $R=50$ each, all of which must be powered.
\item \textbf{Quality control:} both-branch collapse, reported outside the gate.
\end{itemize}
In advance, at null rejection rate $0.05$ and power $0.3$ in each alternative, the gate passes with probability $\numOCvThreeGate$.

\begin{table}[h]
\centering
\caption{Contract of the calibration cells (closed form, standardized scale). ``Changed'' compares the sequential with the direct branch; ``truth'' compares the sequential branch with the true conditional. $\E\,\Delta$CRPS is the expected CRPS excess over the true conditional.}
\label{tab:contract}
\scriptsize
\setlength{\tabcolsep}{2pt}
% AUTO-GENERATED by scripts/make_paper_assets.py from committed raw results. Do not edit by hand.
\begin{tabular}{lcccrr}
\toprule
Condition & target law & proj.\ joint law & truth (target) & sd ratio & $\E\,\Delta$CRPS$\times10^3$ \\
 & \multicolumn{2}{c}{changed between branches?} & seq.\ vs.\ truth & & seq.\ / direct \\
\midrule
Exact sequential & no & no & no & 1.000 & 0.00 / 0.00 \\
Chain $z\mid y$ (Markov) & no & no & no & 1.000 & 0.00 / 0.00 \\
Corr.\ scale $-1$ & no & yes & no & 1.000 & 0.00 / 0.00 \\
Mean shift $0.1\tau$ & yes & yes & yes & 1.000 & 2.81 / 0.00 \\
Variance ratio 1.25 & yes & yes & yes & 1.118 & 1.84 / 0.00 \\
Collapse, intermediate only & yes & yes & yes & 0.735 & 11.32 / 0.00 \\
Collapse, both branches & no & no & yes & 0.000 & 232.52 / 232.52 \\
\bottomrule
\end{tabular}

\end{table}

\paragraph{Calibration result (single run).}
The v3 calibration ran once, at the pre-registered commit, in \numCalibWall{} minutes of wall time on two processes (\numCalibCPU{} CPU-minutes). \Cref{tab:calib} lists every cell.
The gate outcome was \textbf{\numCalibGate}: primary tolerance \numCalibGOne, null-family excess check \numCalibGTwo, power family \numCalibGThree.
\begin{itemize}
\item \textbf{Null validity.} The primary null rejected \numCalibPrimK{} times (one-sided 95\% lower and upper bounds \numCalibPrimLo{} and \numCalibPrimHi, an equal-tailed two-sided 90\% interval; two-sided 95\% CP [\numCalibPrimTwoLo, \numCalibPrimTwoHi]). The joint-level exact null, the target-level correlation null and the Markov chain null rejected \numCalibNullJ, \numCalibNullCorrT{} and \numCalibNullChain.
\item \textbf{Power.} All four fixed alternatives were powered: $0.1\tau$ shift \numCalibPowShift, variance ratio $1.25$ \numCalibPowVar, one-sided collapse \numCalibPowCollInt, joint correlation flip \numCalibPowCorrJ.
\item \textbf{Quality control.} The compatibility test did not reject both-branch collapse in any replicate, as it should not. The quality metrics flagged every replicate: CRPS excess over the truth about $0.235$, sd ratio $0$ (\cref{tab:calibq}).
\item \textbf{Negative controls.} The always-accept and always-reject rules fail the gate.
\end{itemize}

This is an engineering result on the Gaussian probe at this design.
It shows no excess rejection at the tested nulls and adequate power against the tested alternatives.
It does not establish level $0.05$; the level guarantee rests on the permutation argument under \cref{ass:design}.
It does not show that the checker is calibrated for learned generators, whose samples need not satisfy the same assumptions.
It also leaves the v1 record unchanged: the v1 gate failed, and the v1 checker was never cleared.


\paragraph{Consistency and quality answer different questions.}
Per-replicate CRPS excess over the true conditional is reported in \cref{tab:calibq} (\cref{app:calibq}).
Its per-replicate interval excludes $0$ in 12\% of replicates for the $0.1\tau$ shift, and in 14\% for the variance alternative. Test~P rejected 90\% and 52\% of those replicates.
The two are detection rates for \emph{different} nulls, so neither number is a power comparison for the other's hypothesis.
\begin{itemize}
\item The CRPS interval asks whether the sequential branch is worse than the truth, using one held-out observation per example.
\item Test~P asks whether the two branches have the same law.
\end{itemize}
Both-branch collapse shows the converse case: it passes Test~P and fails the quality check in every replicate.



\section{Mixture Probe and Learner Design}
\label{app:mixture}
\paragraph{Source.}
The source is a mixture of three components with weights $(0.5,0.3,0.2)$ and means $(-1,-1,0)$, $(1,1.5,1)$ and $(0,0.5,-1.5)$.
\begin{itemize}
\item Component standard deviations: $(0.6,0.8,0.5)$, $(0.5,0.6,0.7)$ and $(0.8,0.5,0.6)$.
\item Component correlation matrices, as $(\rho_{xy},\rho_{xz},\rho_{yz})$: $(0.5,0.2,0.4)$, $(-0.3,0.4,0.2)$ and $(0.2,-0.5,0.3)$.
\end{itemize}
All three covariances were checked positive definite before any data were drawn, and the parameters were not tuned.
\paragraph{Splits and seeds.}
The splits are $2000$ train, $200$ dev and $200$ test draws. Data, model and sampling seeds use separate hash keys.
\paragraph{Fitted control.}
EM with $K=2$, $500$ maximum iterations, tolerance $10^{-8}$ and covariance regularization $10^{-6}$. It is initialized by $k$-means++ seeding with the model seed.
All marginals and conditionals are computed from the fitted density.
\paragraph{Learners (not run).}
\begin{itemize}
\item \textbf{Loss:} linear-path conditional flow matching.
\item \textbf{Independent arm:} four MLPs (width 64, depth 3), one per pattern $z\mid x$, $y\mid x$, $z\mid x,y$ and $y,z\mid x$, with $5000$ updates each.
\item \textbf{Shared arm:} one MLP (width 128, depth 3) over $(x,y,z)$ whose input marks each coordinate as observed, target or absent, with $20000$ updates and uniform pattern sampling.
\item \textbf{Budgets:} both arms get $20000$ updates in total ($4\times5000$ for the independent arm) and equal expected pattern exposure. Parameter counts are close ($34757$ and $34819$), but each shared update runs a network about four times larger, so the arms are not equal-compute.
\item \textbf{Optimizer:} Adam, learning rate $10^{-3}$, batch size $256$.
\item \textbf{Solver:} Euler from $\mathcal N(0,I)$ with CFG off. The fine level (128) is primary, fixed in advance; the coarse level (32) is a sensitivity check. An arm is labeled inconclusive if the two levels' gap estimates differ by more than $1.645$ combined standard errors, or if its dev loss has not plateaued. Agreement does not prove convergence.
\item \textbf{Checkpoint and comparison:} the final iterate, with no selection. The primary comparison is the paired per-example difference of the unbiased gap estimates of the two arms, over the same test examples.
\item \textbf{Pre-execution fixes:} the shared input gained a target mask, and each independent network got its own optimizer.
\item \textbf{Blocker:} the adapter is written but untested, because no tensor library is available.
\end{itemize}

\section{Pre-Registration Record and Deviations}
\label{app:prereg}
\begin{itemize}
\item \textbf{v1.} Hypotheses EH1--EH6, the statistic, $\alpha=0.05$ and the stop rule were fixed before the test run.
  \begin{itemize}
  \item EH2--EH6 matched their predictions. EH1 failed and is kept as failed.
  \item Post hoc we found: the mis-specified EH1 criterion; the seed flaw that shares noise across joints; and a rendering bug in a table script, which did not affect the raw data.
  \item A reported-statistics error in our own audit notes (a miscounted number of cells) was corrected in the record.
  \end{itemize}
\item \textbf{P5-E8-EXACT} was fixed before running.
  \begin{itemize}
  \item Components A and B passed their deterministic criteria.
  \item In component C, the three pre-registered null cells showed no excess, and the pre-registered powered alternative was powered. The three ``report'' alternatives were not powered.
  \item A development smoke run with a different seed preceded it.
  \end{itemize}
\item \textbf{P5-E8-CALIB, v2} was fixed before any E8 result, except its resource estimate, which it declared would be filled from the E8 timing smoke.
  \begin{itemize}
  \item A pre-run contract audit showed that its gate would fail an exactly valid test with probability $1-\numOCvTwo$.
  \item The audit also found that the E8 seed key omitted the test level, and that v2 had no quality-control family.
  \item v2 was never run and is preserved.
  \end{itemize}
\item \textbf{P5-E8-CALIB, v3} replaced v2 before any calibration result.
  \begin{itemize}
  \item It has three families: null validity, power, and quality control.
  \item Primary null: $R=400$ with a tolerance criterion. The other null cells use Bonferroni excess flags.
  \item Seeds include the test level.
  \item The gate's operating characteristics were computed in advance: P(pass $\mid$ valid, power $0.3$) $=\numOCvThreeGate$.
  \item A timing smoke with a different seed preceded the single run and is not reported.
  \end{itemize}
\item \textbf{Exploratory analyses.} The re-analysis in \cref{app:reanalysis,app:znull} and the closed-form curves in \cref{fig:numerics} were computed after the v1 results were seen.
\end{itemize}

\section{Errata to Earlier Versions}
\label{app:errata}
\begin{enumerate}
\item \textbf{Size.} ``Size exactly $0.05$ for $B=199$'' is replaced by ``level $\le\alpha$''. Exact size needs, in addition, that the permuted statistics have no ties almost surely, which is not guaranteed (\cref{sec:tests}).
\item \textbf{Intervals.} Pairs of one-sided 95\% bounds were written as if they were one interval. Such a pair is an equal-tailed two-sided 90\% interval; two-sided 95\% intervals are now reported.
\item \textbf{CRPS vs.\ Test~P.} Their detection rates concern different nulls and are not a power comparison.
\item \textbf{Scope of the calibration pass.} The Gaussian calibration pass is engineering evidence for one implementation and design. It is not an empirical guarantee on new data.
\item \textbf{Naming.} A network shared across conditionals is a shared conditional network, not a coherent joint.
\end{enumerate}

\section{Compute}
\label{app:compute}
All runs are single-process pure Python on CPU.
\begin{itemize}
\item v1 test run: \numRuntimeVone\,s.
\item P5-E8-EXACT: \numRuntimeEeight\,s.
\item Development, timing and analysis runs are listed in the supplementary compute ledger.
\item P5-E8-CALIB v3: \numCalibWall{} minutes of wall time on two processes (\numCalibCPU{} CPU-minutes); 920 tasks, run once.
\item A calibration timing smoke with a different seed is listed in the ledger and not reported.
\end{itemize}

\section{Reproducibility}
\label{app:repro}
The supplementary code regenerates every number in this paper:
\begin{itemize}
\item \texttt{run\_first\_run.py --split test} reproduces the v1 raw results.
\item \texttt{run\_e8\_exact.py} reproduces P5-E8-EXACT.
\item \texttt{reanalyse\_v1.py} reproduces the exploratory re-analysis.
\item \texttt{make\_paper\_assets.py} regenerates the tables, figure data and numeric macros from committed raw results.
\end{itemize}
Unit tests cover the linear algebra, the conditionals and perturbations, the sorted energy distance against brute force, the fair CRPS against its Gaussian closed form, the affine PF-ODE map against the sampler, the V/U equivalence and exact-enumeration super-uniformity of Test~P, and the Clopper--Pearson bounds. They are numerical checks, not proofs.

\end{document}
